% Auto-split to keep each TeX source < 800 lines.

\paragraph{纤维谱的两态一致渐近与分位阈值（bin-fold）}
本段在命题 \ref{prop:fold-bin-fiber-tail} 的尾部表示之上，给出二进制区间折叠
\(\mathrm{Fold}^{\mathrm{bin}}_m:\{0,\dots,2^m-1\}\to X_m\)
的纤维多重度谱在 \(m\to\infty\) 下的一个闭式极限图景：在指数主阶上，纤维规模退化为仅由末位 \(w_m\) 决定的两态常数；由此导出与 prime-register 分位阈值同型的显式指数常数 \(\log(2/\varphi)\)。

\begin{theorem}[纤维大小的两态一致渐近（bin-fold）]\label{thm:fold-bin-two-state-asymptotic}
\leanverified{paper\_fold\_bin\_two\_state\_asymptotic}
存在绝对常数 \(C>0\)，使得对所有充分大的 \(m\) 以及任意 \(w=(w_1,\dots,w_m)\in X_m\)，二进制区间折叠的纤维多重度满足
\begin{equation}\label{eq:fold-bin-two-state-asymptotic}
\boxed{\;
\Bigl|\,d^{\mathrm{bin}}_m(w)\ -\ \varphi^{-w_m}\Bigl(\frac{2}{\varphi}\Bigr)^m\,\Bigr|\ \le\ C\;}
\end{equation}
其中 \(\varphi=\frac{1+\sqrt5}{2}\)。
等价地，
\[
\Bigl(\frac{\varphi}{2}\Bigr)^m d^{\mathrm{bin}}_m(w)
=\varphi^{-w_m}+O\!\Bigl(\Bigl(\frac{\varphi}{2}\Bigr)^m\Bigr),
\qquad (m\to\infty),
\]
且误差在 \(w\in X_m\) 上一致。
\end{theorem}

\begin{proof}
由命题 \ref{prop:fold-bin-fiber-tail}，
\(
d^{\mathrm{bin}}_m(w)
\)
等价于计数所有满足相邻约束与边界约束的尾部位串 \((c_{m+1},\dots,c_K)\)，并同时满足阈值约束
\[
V_m(w)+\sum_{k=m+1}^{K}c_k F_{k+1}\le 2^m-1,
\]
其中 \(K=K(m)\)（定义 \ref{def:K-of-m}）。
记 \(L:=K-m\ge 0\)，并令
\(
d_i:=c_{m+i}\ (1\le i\le L)
\)。
定义两组尾部线性函数
\[
x(d):=\sum_{i=1}^{L} d_i F_{i+1},
\qquad
y(d):=\sum_{i=1}^{L} d_i F_i.
\]
\par
\noindent\textbf{(1) 尾部和的 Fibonacci 分解。}
由恒等式
\(
F_{m+i+1}=F_{m+1}F_{i+1}+F_m F_i
\)
（对 \(i\) 归纳即可），得到
\begin{equation}\label{eq:fold-bin-tail-sum-split}
\sum_{i=1}^{L} d_i F_{m+i+1}=F_{m+1}x(d)+F_m y(d).
\end{equation}
\par
\noindent\textbf{(2) \(x(d)\) 与 \(y(d)\) 的有界偏差。}
由 Binet 公式推出的精确恒等式
\(
F_{i+1}-\varphi F_i=(-1/\varphi)^i
\)
得
\[
x(d)-\varphi y(d)=\sum_{i=1}^{L} d_i\Bigl(-\frac{1}{\varphi}\Bigr)^i=:S(d).
\]
由于可容许性 \(d_id_{i+1}=0\) 禁止相邻指标同时取 \(1\)，故
\[
-1\ \le\ S(d)\ \le\ \sum_{k\ge 1}\varphi^{-2k}=\frac{1}{\varphi},
\]
从而
\begin{equation}\label{eq:fold-bin-y-as-x-over-phi}
y(d)=\frac{x(d)}{\varphi}+O(1),
\end{equation}
且 \(O(1)\) 绝对有界并与 \(m,w\) 无关。
\par
\noindent\textbf{(3) 阈值约束化归为对 \(x(d)\) 的截断计数。}
将 \eqref{eq:fold-bin-y-as-x-over-phi} 代入 \eqref{eq:fold-bin-tail-sum-split}，并使用
\(
F_{m+1}=\varphi F_m+(-1/\varphi)^m
\)
得到
\[
\sum_{i=1}^{L} d_i F_{m+i+1}
=F_m\Bigl(\Bigl(\varphi+\frac{1}{\varphi}\Bigr)x(d)+O(1)\Bigr)+O(1)
=\sqrt5\,F_m\,x(d)+O(F_m)+O(1).
\]
记 \(B_w:=2^m-1-V_m(w)\)。
由 \(0\le V_m(w)\le F_{m+2}-1\) 与 \(F_{m+2}=\Theta(\varphi^m)\) 得
\begin{equation}\label{eq:fold-bin-Bw-as-2m}
B_w=2^m+O(\varphi^m).
\end{equation}
因此存在绝对常数使
\begin{equation}\label{eq:fold-bin-d-count-x-threshold}
d^{\mathrm{bin}}_m(w)
=\#\Bigl\{d:\ d\ \text{可容许},\ w_m d_1=0,\ x(d)\le \frac{B_w}{\sqrt5\,F_m}+O(1)\Bigr\}.
\end{equation}
\par
\noindent\textbf{(4) 主尺度与末位两态因子。}
由 Binet 渐近 \(F_m=\varphi^m/\sqrt5+O(\varphi^{-m})\) 与 \eqref{eq:fold-bin-Bw-as-2m} 得
\begin{equation}\label{eq:fold-bin-x-threshold-main}
\frac{B_w}{\sqrt5\,F_m}=\Bigl(\frac{2}{\varphi}\Bigr)^m+O(1).
\end{equation}
另一方面，长度 \(L\) 的可容许尾部字数为 \(F_{L+2}\)；
而首位约束 \(d_1=0\) 的可容许字数为 \(F_{L+1}\)，故其比值趋于 \(1/\varphi\)。
将该两态比值与 \eqref{eq:fold-bin-d-count-x-threshold}--\eqref{eq:fold-bin-x-threshold-main} 合并，得到
\[
d^{\mathrm{bin}}_m(w)
=\Bigl(\frac{2}{\varphi}\Bigr)^m\times
\begin{cases}
1+O\bigl((\varphi/2)^m\bigr),& w_m=0,\\[2pt]
\varphi^{-1}+O\bigl((\varphi/2)^m\bigr),& w_m=1,
\end{cases}
\]
从而得到 \eqref{eq:fold-bin-two-state-asymptotic}。证毕。
\end{proof}

\begin{corollary}[两点极限律（bin-fold）]\label{cor:fold-bin-two-point-limit-law}
\leanverified{paper\_fold\_bin\_two\_point\_limit\_law}
令 \(U_m\sim\mathrm{Unif}\{0,\dots,2^m-1\}\)，并令 \(W_m:=\mathrm{Fold}^{\mathrm{bin}}_m(U_m)\in X_m\)。
定义推前分布
\[
\mu_m(w):=\mathbb{P}(W_m=w)=\frac{d^{\mathrm{bin}}_m(w)}{2^m},\qquad w\in X_m,
\]
以及尺度化随机变量
\(
Z_m:=\bigl(\frac{\varphi}{2}\bigr)^m d^{\mathrm{bin}}_m(W_m)
\)。
则 \(Z_m\) 在分布上收敛到两点分布 \(Z\)：
\[
Z=
\begin{cases}
1,& \text{以概率 }\displaystyle \frac{\varphi}{\sqrt5},\\[6pt]
\varphi^{-1},& \text{以概率 }\displaystyle \frac{1}{\varphi\sqrt5}.
\end{cases}
\]
\end{corollary}

\begin{proof}
由定理 \ref{thm:fold-bin-two-state-asymptotic}，有 \(Z_m=\varphi^{-W_m(m)}+o(1)\)（在概率意义下），故只需计算末位 \(W_m(m)=w_m\) 的极限权重。
注意 \(\card{\{w\in X_m:w_m=0\}}=F_{m+1}\)、\(\card{\{w\in X_m:w_m=1\}}=F_m\)。
并由定理主项 \(d^{\mathrm{bin}}_m(w)\sim \varphi^{-w_m}(2/\varphi)^m\) 得
\[
\mathbb{P}(w_m=0)
=\sum_{w_m=0}\mu_m(w)
\sim
F_{m+1}\cdot \frac{(2/\varphi)^m}{2^m}
=\frac{F_{m+1}}{\varphi^m}\to \frac{\varphi}{\sqrt5},
\]
同理 \(\mathbb{P}(w_m=1)\to \frac{1}{\varphi\sqrt5}\)。证毕。
\end{proof}

\begin{corollary}[分位阈值的显式指数常数（bin-fold）]\label{cor:fold-bin-quantile-threshold-constant}
在与定义 \ref{def:pom-gamma-quantile} 同型的分位口径下，对 bin-fold 定义纤维指数
\[
\gamma^{\mathrm{bin}}_m(w):=\frac{1}{m}\log d^{\mathrm{bin}}_m(w),\qquad w\in X_m,
\]
并令 \(\gamma^{\mathrm{bin}}_m(\varepsilon)\) 为其在输出分布 \(\mu_m(w)=d^{\mathrm{bin}}_m(w)/2^m\) 下的尾部分位，
\[
\gamma^{\mathrm{bin}}_{m}(\varepsilon):=\inf\Bigl\{\gamma:\ \mathbb{P}_{\mu_m}\bigl(\gamma^{\mathrm{bin}}_m(W_m)>\gamma\bigr)\le \varepsilon\Bigr\},
\qquad
B^{\mathrm{bin}}_m(\varepsilon):=\exp\!\bigl(m\,\gamma^{\mathrm{bin}}_m(\varepsilon)\bigr).
\]
则当 \(m\to\infty\)，有分段常数极限
\[
\Bigl(\frac{\varphi}{2}\Bigr)^m B^{\mathrm{bin}}_m(\varepsilon)\ \longrightarrow\
\begin{cases}
1,& 0<\varepsilon<\displaystyle\frac{\varphi}{\sqrt5},\\[6pt]
\varphi^{-1},& \displaystyle\frac{\varphi}{\sqrt5}<\varepsilon<1.
\end{cases}
\]
\leanverified{paper\_fold\_bin\_quantile\_threshold\_constant}
特别地，对任意固定 \(\varepsilon\in(0,\varphi/\sqrt5)\)，有
\[
B^{\mathrm{bin}}_m(\varepsilon)
=\Bigl(\frac{2}{\varphi}\Bigr)^m\,(1+o(1)),
\qquad
\frac{1}{m}\log B^{\mathrm{bin}}_m(\varepsilon)\to \log\frac{2}{\varphi}.
\]
\end{corollary}

\begin{proof}
由推论 \ref{cor:fold-bin-two-point-limit-law}，在概率意义下有
\(
(\varphi/2)^m d^{\mathrm{bin}}_m(W_m)\Rightarrow Z\in\{1,\varphi^{-1}\}
\)；
分位极限因此退化为两点分布的分位函数，得到所示分段极限。
最后一式由两端同取 \(\log\) 并除以 \(m\) 得到。证毕。
\end{proof}

\paragraph{输出分布的常数级非均匀性：KL 极限、幂和常数与 R\'enyi 速率塌缩}
两态一致渐近不仅给出分位阈值的显式指数常数，还允许把输出分布 \(\mu_m\) 相对于稳定层均匀分布的偏离压缩为\emph{常数级}账本：
在指数尺度上 \(H_q(\mu_m)/m\) 对一切 \(q>0\) 均塌缩到 \(\log\varphi\)，而所有非均匀性只残留在 \(O(1)\) 的漂移项中。

\begin{proposition}[bin-fold 输出非均匀性的极限常数律：\texorpdfstring{$D_{\mathrm{KL}}(\mu_m\Vert u_m)\to C_{\mathrm{KL}}$}{DKL(mu||u)->CKL}]\label{prop:fold-bin-kl-constant}
\leanverified{paper\_fold\_bin\_kl\_constant}
沿用推论 \ref{cor:fold-bin-two-point-limit-law} 的记号。
令 \(u_m\) 为 \(X_m\) 上的均匀分布：\(u_m(w)=1/|X_m|\)。
定义相对熵（KL 散度）
\[
D_{\mathrm{KL}}(\mu_m\Vert u_m)
:=\sum_{w\in X_m}\mu_m(w)\log\frac{\mu_m(w)}{u_m(w)}.
\]
则存在极限常数
\[
C_{\mathrm{KL}}
=2\log\varphi-\frac12\log 5-\frac{\log\varphi}{\varphi\sqrt5},
\]
使得
\[
D_{\mathrm{KL}}(\mu_m\Vert u_m)\ \longrightarrow\ C_{\mathrm{KL}}
\qquad(m\to\infty).
\]
等价地，纤维信息损失的期望满足
\[
\mathbb{E}_{\mu_m}\bigl[\log d^{\mathrm{bin}}_m(W_m)\bigr]
=m\log\!\Bigl(\frac{2}{\varphi}\Bigr)-\frac{\log\varphi}{\varphi\sqrt5}+o(1),
\]
并且输出 Shannon 熵呈现常数级熵缺口
\[
H(\mu_m)=\log|X_m|-C_{\mathrm{KL}}+o(1),
\qquad
H(\mu_m):=-\sum_{w\in X_m}\mu_m(w)\log \mu_m(w).
\]
\end{proposition}

\begin{proof}
由 \(\mu_m(w)=d^{\mathrm{bin}}_m(w)/2^m\) 与 \(u_m(w)=1/|X_m|\) 得恒等式
\begin{equation}\label{eq:fold-bin-kl-identity}
D_{\mathrm{KL}}(\mu_m\Vert u_m)
=\mathbb{E}_{\mu_m}\bigl[\log d^{\mathrm{bin}}_m(W_m)\bigr]-m\log 2+\log|X_m|.
\end{equation}
由定理 \ref{thm:fold-bin-two-state-asymptotic} 的一致主项，
\[
\log d^{\mathrm{bin}}_m(w)
=m\log\!\Bigl(\frac{2}{\varphi}\Bigr)-w_m\log\varphi+O\!\Bigl(\Bigl(\frac{\varphi}{2}\Bigr)^m\Bigr),
\qquad (w\in X_m),
\]
故
\[
\mathbb{E}_{\mu_m}\bigl[\log d^{\mathrm{bin}}_m(W_m)\bigr]
=m\log\!\Bigl(\frac{2}{\varphi}\Bigr)-\log\varphi\cdot \mathbb{P}(W_m(m)=1)+o(1).
\]
由推论 \ref{cor:fold-bin-two-point-limit-law} 得 \(\mathbb{P}(W_m(m)=1)\to \frac{1}{\varphi\sqrt5}\)，从而得到所示期望渐近。
另一方面，由 Binet 渐近 \(F_{m+2}=\varphi^{m+2}/\sqrt5+o(\varphi^m)\) 得
\[
\log|X_m|=\log F_{m+2}=(m+2)\log\varphi-\frac12\log 5+o(1).
\]
将两式代回 \eqref{eq:fold-bin-kl-identity} 即得 \(D_{\mathrm{KL}}(\mu_m\Vert u_m)\to C_{\mathrm{KL}}\)。
最后，恒等式 \(D_{\mathrm{KL}}(\mu_m\Vert u_m)=\log|X_m|-H(\mu_m)\) 给出熵缺口表达。证毕。
\end{proof}

\begin{proposition}[bin-fold 的全阶碰撞矩主项闭式与指数增长率]\label{prop:fold-bin-power-sum-asymptotic}
对任意固定 \(q\ge 0\)，定义 \(q\) 阶碰撞矩（幂和）
\[
S_q^{\mathrm{bin}}(m):=\sum_{w\in X_m}\bigl(d^{\mathrm{bin}}_m(w)\bigr)^q.
\]
则当 \(m\to\infty\) 时有一致渐近
\[
S_q^{\mathrm{bin}}(m)=\Bigl(\frac{2}{\varphi}\Bigr)^{qm}\Bigl(F_{m+1}+\varphi^{-q}F_m\Bigr)\Bigl(1+O\bigl(q(\varphi/2)^m\bigr)\Bigr),
\]
其中误差在 \(q\) 固定时对 \(m\) 一致。
因此指数增长率极限存在并为
\[
\tau_{\mathrm{bin}}(q):=\lim_{m\to\infty}\frac{1}{m}\log S_q^{\mathrm{bin}}(m)
=q\log\!\Bigl(\frac{2}{\varphi}\Bigr)+\log\varphi,
\]
并且对应的主指数根（Perron 根）为
\[
r_q:=\exp\bigl(\tau_{\mathrm{bin}}(q)\bigr)=\frac{2^q}{\varphi^{q-1}}=2^q\varphi^{1-q}.
\]
\leanpartial{paper\_fold\_bin\_power\_sum\_asymptotic}{the current concrete Lean `momentSum` model has a counterexample at `q = 6`, `m = 8`, so this asymptotic bound is false as stated in the repository definitions}
\end{proposition}

\begin{proof}
由定理 \ref{thm:fold-bin-two-state-asymptotic}，对所有 \(w\in X_m\) 有一致展开
\[
d^{\mathrm{bin}}_m(w)=\varphi^{-w_m}\Bigl(\frac{2}{\varphi}\Bigr)^m\Bigl(1+O\bigl((\varphi/2)^m\bigr)\Bigr),
\]
其中误差项在 \(w\) 上一致。从而对任意固定 \(q\ge 0\) 可得
\[
\bigl(d^{\mathrm{bin}}_m(w)\bigr)^q
=\Bigl(\frac{2}{\varphi}\Bigr)^{qm}\varphi^{-q w_m}\Bigl(1+O\bigl(q(\varphi/2)^m\bigr)\Bigr),
\]
其中误差项在 \(w\in X_m\) 上一致。
将 \(X_m\) 按末位 \(w_m\in\{0,1\}\) 分拆，并用
\(
\card\{w\in X_m:w_m=0\}=F_{m+1}
\)、\(
\card\{w\in X_m:w_m=1\}=F_{m}
\)，得
\[
S_q^{\mathrm{bin}}(m)
=\Bigl(\frac{2}{\varphi}\Bigr)^{qm}\Bigl(F_{m+1}+\varphi^{-q}F_m\Bigr)\Bigl(1+O\bigl(q(\varphi/2)^m\bigr)\Bigr).
\]
两端取对数并除以 \(m\)，再令 \(m\to\infty\)，由 \(F_{m+1}+\varphi^{-q}F_m=\Theta(\varphi^m)\) 得
\(\tau_{\mathrm{bin}}(q)=q\log(2/\varphi)+\log\varphi\)，从而 \(r_q=\exp(\tau_{\mathrm{bin}}(q))=2^q\varphi^{1-q}\)。证毕。
\end{proof}

\begin{theorem}[R\'enyi 速率塌缩与“伪多重分形”消失：\texorpdfstring{$H_q(\mu_m)/m\to\log\varphi$}{Hq(mu)/m->log phi}]\label{thm:fold-bin-renyi-rate-collapse}
沿用推论 \ref{cor:fold-bin-two-point-limit-law} 的 \(\mu_m\) 记号。
对 \(q>0\)、\(q\neq 1\) 定义 R\'enyi 熵
\[
H_q(\mu_m):=\frac{1}{1-q}\log\sum_{w\in X_m}\mu_m(w)^q,
\]
并以极限定义 Shannon 熵 \(H_1(\mu_m):=H(\mu_m)\)。
则对所有 \(q>0\) 都成立统一极限
\[
\lim_{m\to\infty}\frac{1}{m}H_q(\mu_m)=\log\varphi,
\qquad
\lim_{m\to\infty}\frac{1}{m}H_1(\mu_m)=\log\varphi.
\]
进一步，令
\(
\gamma_m:=\frac{1}{m}\log d^{\mathrm{bin}}_m(W_m)
\)（\(W_m\sim\mu_m\)），
则 \(\gamma_m\to\log(2/\varphi)\)，且偏离仅为 \(O(m^{-1})\)。
\end{theorem}

\begin{proof}
当 \(q\neq 1\) 时，由 \(\mu_m(w)=d^{\mathrm{bin}}_m(w)/2^m\) 得
\[
\sum_{w\in X_m}\mu_m(w)^q
=2^{-qm}\,S_q^{\mathrm{bin}}(m).
\]
将命题 \ref{prop:fold-bin-power-sum-asymptotic} 的渐近代入得
\[
\log\sum_{w}\mu_m(w)^q
=-qm\log 2 + m\log r_q + O(1)
=(1-q)m\log\varphi+O(1),
\]
其中用到 \(r_q=2^q/\varphi^{q-1}\)。
两端同除以 \(m(1-q)\) 并令 \(m\to\infty\) 即得 \(\frac{1}{m}H_q(\mu_m)\to\log\varphi\)。
\par
当 \(q=1\) 时，由命题 \ref{prop:fold-bin-kl-constant} 得
\(
H(\mu_m)=\log|X_m|-D_{\mathrm{KL}}(\mu_m\Vert u_m)
=m\log\varphi+O(1)
\)，
从而 \(\frac{1}{m}H_1(\mu_m)\to\log\varphi\)。
\par
最后，由定理 \ref{thm:fold-bin-two-state-asymptotic} 的一致展开
\(
\log d^{\mathrm{bin}}_m(W_m)=m\log(2/\varphi)-W_m(m)\log\varphi+o(1)
\)
且 \(W_m(m)\in\{0,1\}\)，两端除以 \(m\) 得
\(
\gamma_m=\log(2/\varphi)+O(m^{-1})
\)，故结论成立。证毕。
\leanverified{paper\_fold\_bin\_renyi\_rate\_collapse}
\end{proof}

% Auto-split to keep each TeX source < 600 lines.

\paragraph{bin-fold 的可解信息几何：escort 温度族、散度谱与规范密度的两尺度修正}
定理 \ref{thm:fold-bin-two-state-asymptotic} 的两态一致渐近表明：在指数主阶上，纤维多重度仅依赖于末位 \(w_m\)。
这一退化把若干常用的信息几何泛函（escort 温度族、Rényi/Csiszár 散度族以及规范体积密度）统一压缩为同一组显式常数，并给出一致的指数收敛率。

\begin{proposition}[bin-fold 推前分布的末位充分化：全变差塌缩与指数倾斜模型]\label{prop:fold-bin-lastbit-sufficient-tv}
沿用推论 \ref{cor:fold-bin-two-point-limit-law} 的 \(\mu_m\) 记号，并令 \(\ell(w):=w_m\in\{0,1\}\) 为末位指示。
将 \(X_m\) 按末位分成两块
\[
A_i:=\{w\in X_m:\ \ell(w)=i\}\qquad(i\in\{0,1\}),
\]
并定义末位条件化近似（块内均匀化）
\[
\overline\mu_m(w):=\frac{\mu_m(A_{\ell(w)})}{|A_{\ell(w)}|},\qquad w\in X_m.
\]
则存在绝对常数 \(C>0\) 使对所有充分大的 \(m\) 成立
\[
D_{\mathrm{TV}}(\mu_m,\overline\mu_m)
:=\frac12\sum_{w\in X_m}\bigl|\mu_m(w)-\overline\mu_m(w)\bigr|
\ \le\ C\Bigl(\frac{\varphi}{2}\Bigr)^m.
\]
此外，存在 \(\beta_m\in\RR\) 使 \(\overline\mu_m\) 可写为稳定层均匀分布 \(u_m(w)=1/|X_m|\) 的指数倾斜：
\[
\overline\mu_m(w)=\frac{1}{Z_m}\,u_m(w)\,e^{-\beta_m\ell(w)},
\qquad
Z_m:=\sum_{w\in X_m}u_m(w)e^{-\beta_m\ell(w)},
\]
并且 \(\beta_m\to\log\varphi\)。
\leanverified{paper\_fold\_bin\_lastbit\_sufficient\_tv}
\end{proposition}
\begin{proof}
由定理 \ref{thm:fold-bin-two-state-asymptotic} 的一致界 \eqref{eq:fold-bin-two-state-asymptotic}，在每个末位块 \(A_i\) 内，
\(
d^{\mathrm{bin}}_m(w)
\)
与其块内平均值之差为 \(O(1)\)（一致于 \(w\in A_i\)），从而
\[
\sup_{w\in A_i}\bigl|\mu_m(w)-\overline\mu_m(w)\bigr|
=\sup_{w\in A_i}\Bigl|\frac{d^{\mathrm{bin}}_m(w)}{2^m}-\frac{1}{|A_i|}\sum_{w'\in A_i}\frac{d^{\mathrm{bin}}_m(w')}{2^m}\Bigr|
\ \le\ \frac{C_1}{2^m}.
\]
对两块分别求和得
\[
\sum_{w\in X_m}\bigl|\mu_m(w)-\overline\mu_m(w)\bigr|
\le
|X_m|\cdot \frac{C_1}{2^m}
=O\!\Bigl(\Bigl(\frac{\varphi}{2}\Bigr)^m\Bigr),
\]
其中用到 \(|X_m|=F_{m+2}=\Theta(\varphi^m)\)，故得到所示全变差界。
\par
对指数倾斜表示，注意 \(\overline\mu_m\) 在每个 \(A_i\) 上为常数：\(\overline\mu_m|_{A_i}=\mu_m(A_i)/|A_i|\)。
令
\[
e^{-\beta_m}:=\frac{\mu_m(A_1)/|A_1|}{\mu_m(A_0)/|A_0|},
\]
则对任意 \(w\in A_i\) 有 \(\overline\mu_m(w)\propto e^{-\beta_m i}\)，从而存在归一化常数 \(Z_m\) 使
\(
\overline\mu_m(w)=Z_m^{-1}u_m(w)e^{-\beta_m\ell(w)}
\)。
\par
最后，由推论 \ref{cor:fold-bin-two-point-limit-law} 的块质量极限
\(\mu_m(A_0)\to \varphi/\sqrt5\)、\(\mu_m(A_1)\to 1/(\varphi\sqrt5)\)，以及 \(|A_0|=F_{m+1}\)、\(|A_1|=F_m\) 给出的
\(|A_1|/|A_0|\to \varphi^{-1}\)，可得
\[
e^{-\beta_m}
=\frac{\mu_m(A_1)}{\mu_m(A_0)}\cdot\frac{|A_0|}{|A_1|}
\longrightarrow
\frac{\varphi^{-2}}{1}\cdot \varphi
=\varphi^{-1},
\]
因此 \(\beta_m\to\log\varphi\)。证毕。
\end{proof}

\begin{proposition}[escort 温度族的末位极限律与纤维信息常数项（bin-fold）]\label{prop:fold-bin-escort-lastbit}
\leanverified{paper\_fold\_bin\_escort\_lastbit}
对任意 \(q\ge 0\)，定义 escort 分布
\[
\pi_{m,q}(w):=\frac{\bigl(d^{\mathrm{bin}}_m(w)\bigr)^q}{S_q^{\mathrm{bin}}(m)},\qquad w\in X_m,
\]
并令随机变量 \(W^{(q)}_m\sim\pi_{m,q}\)。
设
\[
\theta(q):=\frac{1}{1+\varphi^{q+1}}.
\]
则当 \(m\to\infty\) 时有
\[
\mathbb{P}\bigl(W^{(q)}_m(m)=1\bigr)\to \theta(q),
\qquad
\mathbb{P}\bigl(W^{(q)}_m(m)=0\bigr)\to 1-\theta(q),
\]
并且对固定 \(q\) 收敛速度为 \(O\bigl((\varphi/2)^m\bigr)\)。

进一步，定义 \(q\)-escort 纤维信息
\[
\kappa^{\mathrm{bin}}_m(q):=\mathbb{E}_{\pi_{m,q}}\bigl[\log d^{\mathrm{bin}}_m(W^{(q)}_m)\bigr],
\]
则存在一致展开
\[
\kappa^{\mathrm{bin}}_m(q)=m\log\!\Bigl(\frac{2}{\varphi}\Bigr)-\theta(q)\log\varphi+O\bigl((\varphi/2)^m\bigr)
=m\log\!\Bigl(\frac{2}{\varphi}\Bigr)-\frac{\log\varphi}{1+\varphi^{q+1}}+O\bigl((\varphi/2)^m\bigr),
\]
其中误差在 \(q\) 固定时对 \(m\) 一致。
\end{proposition}

\begin{proof}
由定义，
\[
\mathbb{P}\bigl(W^{(q)}_m(m)=1\bigr)
=\sum_{w_m=1}\pi_{m,q}(w)
=\frac{\sum_{w_m=1}\bigl(d^{\mathrm{bin}}_m(w)\bigr)^q}{S_q^{\mathrm{bin}}(m)}.
\]
由定理 \ref{thm:fold-bin-two-state-asymptotic} 的一致主项可得
\(
\bigl(d^{\mathrm{bin}}_m(w)\bigr)^q=\bigl(\frac{2}{\varphi}\bigr)^{qm}\varphi^{-qw_m}\bigl(1+O(q(\varphi/2)^m)\bigr)
\)
在 \(w\in X_m\) 上一致。
因此
\[
\sum_{w_m=1}\bigl(d^{\mathrm{bin}}_m(w)\bigr)^q
=F_m\Bigl(\frac{2}{\varphi}\Bigr)^{qm}\varphi^{-q}\Bigl(1+O\bigl(q(\varphi/2)^m\bigr)\Bigr),
\]
而命题 \ref{prop:fold-bin-power-sum-asymptotic} 给出
\(
S_q^{\mathrm{bin}}(m)=\bigl(\frac{2}{\varphi}\bigr)^{qm}\bigl(F_{m+1}+\varphi^{-q}F_m\bigr)\bigl(1+O(q(\varphi/2)^m)\bigr)
\)。
两式相除得到
\[
\mathbb{P}\bigl(W^{(q)}_m(m)=1\bigr)
=\frac{F_m\varphi^{-q}}{F_{m+1}+\varphi^{-q}F_m}+O\bigl((\varphi/2)^m\bigr)
\longrightarrow\ \frac{1}{1+\varphi^{q+1}},
\]
其中用到 \(F_{m+1}/F_m\to\varphi\)。
另一端概率由互补得到。

对 \(\kappa^{\mathrm{bin}}_m(q)\)，由定理 \ref{thm:fold-bin-two-state-asymptotic} 的一致对数展开
\[
\log d^{\mathrm{bin}}_m(w)
=m\log\!\Bigl(\frac{2}{\varphi}\Bigr)-w_m\log\varphi+O\bigl((\varphi/2)^m\bigr),
\qquad (w\in X_m),
\]
对 \(\pi_{m,q}\) 取期望并代入 \(\mathbb{P}(W^{(q)}_m(m)=1)=\theta(q)+O((\varphi/2)^m)\) 即得结论。
\end{proof}

\begin{proposition}[escort--escort 的温度间 KL 极限（bin-fold）]\label{prop:fold-bin-escort-escort-kl}
对 \(p,q\ge 0\)，令 \(\pi_{m,q}\) 如命题 \ref{prop:fold-bin-escort-lastbit}。
则极限
\[
\lim_{m\to\infty}D_{\mathrm{KL}}(\pi_{m,q}\Vert \pi_{m,p})
\]
存在，并可写成一维 Bernoulli KL 的闭式。设 \(\theta(\cdot)\) 如命题 \ref{prop:fold-bin-escort-lastbit} 中所定义，则
\[
\lim_{m\to\infty}D_{\mathrm{KL}}(\pi_{m,q}\Vert \pi_{m,p})
=\theta(q)\log\Bigl(\frac{\theta(q)}{\theta(p)}\Bigr)+\bigl(1-\theta(q)\bigr)\log\Bigl(\frac{1-\theta(q)}{1-\theta(p)}\Bigr).
\]
\leanverified{paper\_fold\_bin\_escort\_escort\_kl}
\end{proposition}

\begin{proof}
由定理 \ref{thm:fold-bin-two-state-asymptotic}，对任意固定 \(q\ge 0\) 有
\[
\pi_{m,q}(w)=\frac{\varphi^{-qw_m}}{F_{m+1}+\varphi^{-q}F_m}\Bigl(1+O\bigl((\varphi/2)^m\bigr)\Bigr),
\qquad (w\in X_m),
\]
其中误差在 \(w\) 上一致；同理对 \(p\) 成立。
因此在 \(w_m\) 给定的条件下，\(\pi_{m,q}\) 与 \(\pi_{m,p}\) 在该扇区内均渐近为均匀分布，而比值 \(\pi_{m,q}(w)/\pi_{m,p}(w)\) 亦在扇区内渐近为常数。
将 KL 散度按两扇区 \(w_m\in\{0,1\}\) 分解，并用
\(\mathbb{P}_{\pi_{m,q}}(w_m=1)\to \theta(q)\)、\(\mathbb{P}_{\pi_{m,p}}(w_m=1)\to \theta(p)\)
即可得到所示 Bernoulli KL 极限。
\end{proof}

\begin{corollary}[两尺度极限残差的两点律（bin-fold）]\label{cor:fold-bin-escort-two-scale-residual}
\leanverified{paper\_fold\_bin\_escort\_two\_scale\_residual}
对任意 \(q\ge 0\) 令 \(W^{(q)}_m\sim\pi_{m,q}\)，并定义尺度化残差变量
\[
Z^{(q)}_m:=\Bigl(\frac{\varphi}{2}\Bigr)^m d^{\mathrm{bin}}_m\bigl(W^{(q)}_m\bigr).
\]
则 \(Z^{(q)}_m\Rightarrow Z^{(q)}\) 且 \(Z^{(q)}\) 为两点分布
\[
Z^{(q)}\in\{1,\varphi^{-1}\},
\qquad
\mathbb{P}\bigl(Z^{(q)}=\varphi^{-1}\bigr)=\theta(q),
\quad
\mathbb{P}\bigl(Z^{(q)}=1\bigr)=1-\theta(q).
\]
\end{corollary}

\begin{proof}
由定理 \ref{thm:fold-bin-two-state-asymptotic}，
\(
Z^{(q)}_m=\varphi^{-W^{(q)}_m(m)}+o(1)
\)
（在概率意义下），而 \(W^{(q)}_m(m)\) 的极限分布由命题 \ref{prop:fold-bin-escort-lastbit} 给出，故结论成立。
\end{proof}

\begin{theorem}[escort 温度族的两态 Gibbs 化与末位统计的渐近充分性]\label{thm:killo-fold-bin-escort-onebit-fisher}
对每个 \(m\ge 1\)，记
\[
A_0:=\{w\in X_m:\ w_m=0\},
\qquad
A_1:=\{w\in X_m:\ w_m=1\},
\]
并令 \(u_m\) 为 \(X_m\) 上的均匀分布。
对任意 \(q\ge 0\)，定义 escort 分布
\[
\pi_{m,q}(w):=\frac{\bigl(d_m^{\mathrm{bin}}(w)\bigr)^q}{S_q^{\mathrm{bin}}(m)},
\qquad
S_q^{\mathrm{bin}}(m):=\sum_{x\in X_m}\bigl(d_m^{\mathrm{bin}}(x)\bigr)^q,
\]
以及两态 Gibbs 近似
\[
\nu_{m,q}(w):=\frac{u_m(w)\varphi^{-q w_m}}{Z_{m,q}},
\qquad
Z_{m,q}:=\frac{F_{m+1}+\varphi^{-q}F_m}{F_{m+2}}.
\]
则对任意有界温度区间 \(q\in[0,Q]\)，存在常数 \(C_Q>0\) 使对所有充分大的 \(m\) 成立
\[
\sup_{q\in[0,Q]}\sup_{w\in X_m}
\left|
\frac{\pi_{m,q}(w)}{\nu_{m,q}(w)}-1
\right|
\le C_Q\Bigl(\frac{\varphi}{2}\Bigr)^m.
\]
从而
\[
\sup_{q\in[0,Q]}
\|\pi_{m,q}-\nu_{m,q}\|_{\mathrm{TV}}
=
O_Q\!\Bigl(\Bigl(\frac{\varphi}{2}\Bigr)^m\Bigr).
\]
特别地，末位映射 \(w\mapsto w_m\in\{0,1\}\) 构成 escort 温度族的渐近充分统计量：对 \(i\in\{0,1\}\)，有
\[
\sup_{q\in[0,Q]}
\bigl\|
\pi_{m,q}(\,\cdot\,|A_i)-u_m(\,\cdot\,|A_i)
\bigr\|_{\mathrm{TV}}
=
O_Q\!\Bigl(\Bigl(\frac{\varphi}{2}\Bigr)^m\Bigr).
\]
\end{theorem}
\begin{proof}
由定理 \ref{thm:fold-bin-two-state-asymptotic} 与命题 \ref{prop:fold-bin-power-sum-asymptotic}，对每个有界区间 \(q\in[0,Q]\) 都有一致展开
\[
\bigl(d_m^{\mathrm{bin}}(w)\bigr)^q
=
\Bigl(\frac{2}{\varphi}\Bigr)^{qm}\varphi^{-q w_m}
\Bigl(1+O_Q\!\Bigl(\Bigl(\frac{\varphi}{2}\Bigr)^m\Bigr)\Bigr),
\qquad
(w\in X_m),
\]
以及
\[
S_q^{\mathrm{bin}}(m)
=
\Bigl(\frac{2}{\varphi}\Bigr)^{qm}
\bigl(F_{m+1}+\varphi^{-q}F_m\bigr)
\Bigl(1+O_Q\!\Bigl(\Bigl(\frac{\varphi}{2}\Bigr)^m\Bigr)\Bigr),
\]
其中误差对 \(q\in[0,Q]\) 与 \(w\in X_m\) 一致。
将两式相除，并用 \(u_m(w)=1/F_{m+2}\)，得到
\[
\pi_{m,q}(w)
=
\frac{\varphi^{-q w_m}}{F_{m+1}+\varphi^{-q}F_m}
\Bigl(1+O_Q\!\Bigl(\Bigl(\frac{\varphi}{2}\Bigr)^m\Bigr)\Bigr)
=
\nu_{m,q}(w)
\Bigl(1+O_Q\!\Bigl(\Bigl(\frac{\varphi}{2}\Bigr)^m\Bigr)\Bigr),
\]
即得逐点乘法逼近。

由 \(\sum_{w\in X_m}\nu_{m,q}(w)=1\)，直接得到
\[
\|\pi_{m,q}-\nu_{m,q}\|_{\mathrm{TV}}
\le
\frac12\sum_{w\in X_m}\nu_{m,q}(w)\,
O_Q\!\Bigl(\Bigl(\frac{\varphi}{2}\Bigr)^m\Bigr)
=
O_Q\!\Bigl(\Bigl(\frac{\varphi}{2}\Bigr)^m\Bigr).
\]

再看条件分布。
由定义，\(\nu_{m,q}\) 在每个块 \(A_i\) 上严格均匀；并且
\[
\nu_{m,q}(A_1)=\frac{\varphi^{-q}F_m}{F_{m+1}+\varphi^{-q}F_m},
\qquad
\nu_{m,q}(A_0)=\frac{F_{m+1}}{F_{m+1}+\varphi^{-q}F_m}.
\]
由于 \(q\in[0,Q]\) 且 \(F_{m+1}/F_m\to\varphi\)，存在常数 \(c_Q>0\) 使对所有充分大的 \(m\) 有
\[
c_Q\le \nu_{m,q}(A_i)\le 1-c_Q
\qquad
(i=0,1,\ q\in[0,Q]).
\]
另一方面，由逐点乘法逼近，对 \(w\in A_i\) 有
\[
\pi_{m,q}(w)=\nu_{m,q}(w)\Bigl(1+O_Q\!\Bigl(\Bigl(\frac{\varphi}{2}\Bigr)^m\Bigr)\Bigr),
\qquad
\pi_{m,q}(A_i)=\nu_{m,q}(A_i)\Bigl(1+O_Q\!\Bigl(\Bigl(\frac{\varphi}{2}\Bigr)^m\Bigr)\Bigr).
\]
因此
\[
\pi_{m,q}(w|A_i)
=
\frac{\nu_{m,q}(w)}{\nu_{m,q}(A_i)}
\Bigl(1+O_Q\!\Bigl(\Bigl(\frac{\varphi}{2}\Bigr)^m\Bigr)\Bigr)
=
\nu_{m,q}(w|A_i)
\Bigl(1+O_Q\!\Bigl(\Bigl(\frac{\varphi}{2}\Bigr)^m\Bigr)\Bigr).
\]
而 \(\nu_{m,q}(\cdot|A_i)=u_m(\cdot|A_i)\) 在 \(A_i\) 上均匀，故
\[
\bigl\|
\pi_{m,q}(\,\cdot\,|A_i)-u_m(\,\cdot\,|A_i)
\bigr\|_{\mathrm{TV}}
=
O_Q\!\Bigl(\Bigl(\frac{\varphi}{2}\Bigr)^m\Bigr),
\]
即证。
\end{proof}

\leanverified{paper\_killo\_fold\_bin\_escort\_onebit\_fisher}

\leanverified{paper\_killo\_fold\_bin\_escort\_renyi\_logistic\_geometry}
\begin{theorem}[escort 温度族的全阶 Rényi 信息几何由 logistic 曲线控制]\label{thm:killo-fold-bin-escort-renyi-logistic-geometry}
定义
\[
\theta(q):=\frac{1}{1+\varphi^{q+1}},
\qquad q\ge 0.
\]
则对任意固定 \(p,q\ge 0\) 与任意固定 \(\alpha\in(0,\infty)\setminus\{1\}\)，有
\[
\lim_{m\to\infty}D_\alpha(\pi_{m,q}\Vert \pi_{m,p})
=
\frac{1}{\alpha-1}
\log\!\Bigl(
\theta(q)^\alpha\theta(p)^{1-\alpha}
+
\bigl(1-\theta(q)\bigr)^\alpha\bigl(1-\theta(p)\bigr)^{1-\alpha}
\Bigr).
\]
当 \(\alpha\to 1\) 时，上式连续退化为
\[
\lim_{m\to\infty}D_{\mathrm{KL}}(\pi_{m,q}\Vert \pi_{m,p})
=
\theta(q)\log\frac{\theta(q)}{\theta(p)}
+
\bigl(1-\theta(q)\bigr)\log\frac{1-\theta(q)}{1-\theta(p)}.
\]
此外，函数 \(q\mapsto \theta(q)\) 满足 logistic 微分方程
\[
\theta'(q)=-(\log\varphi)\,\theta(q)\bigl(1-\theta(q)\bigr),
\]
从而其 Fisher 信息为
\[
I(q):=\frac{\theta'(q)^2}{\theta(q)(1-\theta(q))}
=(\log\varphi)^2\theta(q)\bigl(1-\theta(q)\bigr).
\]
并且 \(q\mapsto I(q)\) 在 \([0,\infty)\) 上严格递减。
\end{theorem}
\begin{proof}
记
\[
\theta_m(q):=\frac{\varphi^{-q}F_m}{F_{m+1}+\varphi^{-q}F_m}.
\]
由定理 \ref{thm:killo-fold-bin-escort-onebit-fisher}，对固定 \(p,q\) 有
\[
\pi_{m,q}(w)=\nu_{m,q}(w)\bigl(1+O((\varphi/2)^m)\bigr),
\qquad
\pi_{m,p}(w)=\nu_{m,p}(w)\bigl(1+O((\varphi/2)^m)\bigr),
\]
其中误差在 \(w\in X_m\) 上一致。
因此对固定 \(\alpha\in(0,\infty)\setminus\{1\}\)，有
\[
\pi_{m,q}(w)^\alpha\pi_{m,p}(w)^{1-\alpha}
=
\nu_{m,q}(w)^\alpha\nu_{m,p}(w)^{1-\alpha}
\Bigl(1+O\!\Bigl(\Bigl(\frac{\varphi}{2}\Bigr)^m\Bigr)\Bigr),
\]
并且该误差仍对 \(w\) 一致。
于是
\[
\sum_{w\in X_m}\pi_{m,q}(w)^\alpha\pi_{m,p}(w)^{1-\alpha}
=
\Bigl(\sum_{w\in X_m}\nu_{m,q}(w)^\alpha\nu_{m,p}(w)^{1-\alpha}\Bigr)
\Bigl(1+O\!\Bigl(\Bigl(\frac{\varphi}{2}\Bigr)^m\Bigr)\Bigr).
\]

另一方面，\(\nu_{m,q}\) 在两块上的表达为
\[
\nu_{m,q}(w)=
\begin{cases}
\dfrac{1-\theta_m(q)}{|A_0|}, & w\in A_0,\\[2ex]
\dfrac{\theta_m(q)}{|A_1|}, & w\in A_1.
\end{cases}
\]
于是
\[
\sum_{w\in X_m}\nu_{m,q}(w)^\alpha\nu_{m,p}(w)^{1-\alpha}
=
\bigl(1-\theta_m(q)\bigr)^\alpha\bigl(1-\theta_m(p)\bigr)^{1-\alpha}
+
\theta_m(q)^\alpha\theta_m(p)^{1-\alpha},
\]
其中块大小完全相消。
故
\[
D_\alpha(\nu_{m,q}\Vert \nu_{m,p})
=
\frac{1}{\alpha-1}
\log\!\Bigl(
\theta_m(q)^\alpha\theta_m(p)^{1-\alpha}
+
\bigl(1-\theta_m(q)\bigr)^\alpha\bigl(1-\theta_m(p)\bigr)^{1-\alpha}
\Bigr).
\]
由上一式可知
\[
D_\alpha(\pi_{m,q}\Vert \pi_{m,p})
=
D_\alpha(\nu_{m,q}\Vert \nu_{m,p})+o(1),
\]
再由 \(\theta_m(\cdot)\to \theta(\cdot)\) 即得所示 Rényi 极限。

当 \(\alpha\to 1\) 时，Rényi 散度向 KL 散度连续延拓，得到
\[
\theta(q)\log\frac{\theta(q)}{\theta(p)}
+
\bigl(1-\theta(q)\bigr)\log\frac{1-\theta(q)}{1-\theta(p)},
\]
这与命题 \ref{prop:fold-bin-escort-escort-kl} 一致。

最后，由
\[
\theta(q)=\frac{1}{1+\varphi^{q+1}}
\]
直接求导得
\[
\theta'(q)
=
-\frac{\varphi^{q+1}\log\varphi}{(1+\varphi^{q+1})^2}
=
-(\log\varphi)\,\theta(q)\bigl(1-\theta(q)\bigr).
\]
因此
\[
I(q)=\frac{\theta'(q)^2}{\theta(q)(1-\theta(q))}
=(\log\varphi)^2\theta(q)\bigl(1-\theta(q)\bigr).
\]
又因
\[
0<\theta(q)\le \theta(0)=\frac{1}{1+\varphi}=\varphi^{-2}<\frac12,
\]
且 \(q\mapsto \theta(q)\) 严格递减，而 \(x\mapsto x(1-x)\) 在 \((0,1/2]\) 上严格递增，故
\[
q\mapsto \theta(q)\bigl(1-\theta(q)\bigr)
\]
在 \([0,\infty)\) 上严格递减，从而 \(I(q)\) 亦严格递减。
\end{proof}

\endinput

\begin{theorem}[escort 温度族的全阶 Rényi 熵常数项]\label{thm:killo-fold-bin-escort-renyi-entropy}
\leanverified{paper\_killo\_fold\_bin\_escort\_renyi\_entropy}
沿用定理 \ref{thm:killo-fold-bin-escort-renyi-logistic-geometry} 的 \(\theta(q)\) 记号。
对任意固定 \(q\ge 0\) 与任意固定 \(\alpha\in(0,\infty)\setminus\{1\}\)，escort 分布的 Rényi 熵满足
\[
H_\alpha(\pi_{m,q})
=
m\log\varphi
-\frac12\log 5
+\frac{1}{1-\alpha}
\log\!\Bigl(
\varphi^{1-\alpha}\bigl(1-\theta(q)\bigr)^\alpha+\theta(q)^\alpha
\Bigr)
+o(1).
\]
当 \(\alpha=1\) 时，若记
\[
h_2(t):=-t\log t-(1-t)\log(1-t),
\qquad 0\le t\le 1,
\]
则有
\[
H(\pi_{m,q})
=
m\log\varphi
-\frac12\log 5
+\bigl(1-\theta(q)\bigr)\log\varphi
+h_2(\theta(q))
+o(1).
\]
特别地，对一切固定 \(q\ge 0\) 与一切固定 \(\alpha>0\)，熵率极限满足
\[
\lim_{m\to\infty}\frac{1}{m}H_\alpha(\pi_{m,q})=\log\varphi.
\]
\end{theorem}
\begin{proof}
先对定理 \ref{thm:killo-fold-bin-escort-onebit-fisher} 中的两态 Gibbs 近似 \(\nu_{m,q}\) 计算。
记
\[
\theta_m(q):=\frac{\varphi^{-q}F_m}{F_{m+1}+\varphi^{-q}F_m}.
\]
则
\[
\nu_{m,q}(w)=
\begin{cases}
\dfrac{1-\theta_m(q)}{F_{m+1}}, & w\in A_0,\\[2ex]
\dfrac{\theta_m(q)}{F_m}, & w\in A_1.
\end{cases}
\]
因此当 \(\alpha\neq 1\) 时，
\[
\sum_{w\in X_m}\nu_{m,q}(w)^\alpha
=
F_{m+1}^{1-\alpha}\bigl(1-\theta_m(q)\bigr)^\alpha
+
F_m^{1-\alpha}\theta_m(q)^\alpha.
\]
代回 Rényi 熵定义
\[
H_\alpha(\nu_{m,q})
:=
\frac{1}{1-\alpha}\log\sum_{w\in X_m}\nu_{m,q}(w)^\alpha
\]
并用 Binet 渐近
\[
F_m=\frac{\varphi^m}{\sqrt5}+O(\varphi^{-m}),
\qquad
F_{m+1}=\frac{\varphi^{m+1}}{\sqrt5}+O(\varphi^{-m}),
\]
以及 \(\theta_m(q)\to \theta(q)\)，得到
\[
H_\alpha(\nu_{m,q})
=
m\log\varphi
-\frac12\log 5
+\frac{1}{1-\alpha}
\log\!\Bigl(
\varphi^{1-\alpha}\bigl(1-\theta(q)\bigr)^\alpha+\theta(q)^\alpha
\Bigr)
+o(1).
\]

再由定理 \ref{thm:killo-fold-bin-escort-onebit-fisher}，对固定 \(q\) 有
\[
\sup_{w\in X_m}
\left|
\frac{\pi_{m,q}(w)}{\nu_{m,q}(w)}-1
\right|
=o(1).
\]
因此对固定 \(\alpha>0\)，有一致乘法逼近
\[
\pi_{m,q}(w)^\alpha=\nu_{m,q}(w)^\alpha(1+o(1)),
\qquad (w\in X_m),
\]
从而
\[
\sum_{w\in X_m}\pi_{m,q}(w)^\alpha
=
\Bigl(\sum_{w\in X_m}\nu_{m,q}(w)^\alpha\Bigr)(1+o(1)).
\]
取对数并除以 \(1-\alpha\) 可得
\[
H_\alpha(\pi_{m,q})=H_\alpha(\nu_{m,q})+o(1),
\]
即得 \(\alpha\neq 1\) 的公式。

对 Shannon 熵，\(\nu_{m,q}\) 的块结构给出精确恒等式
\[
H(\nu_{m,q})
=
h_2(\theta_m(q))
+\bigl(1-\theta_m(q)\bigr)\log F_{m+1}
+\theta_m(q)\log F_m.
\]
由 \(\theta_m(q)\to\theta(q)\) 与 Binet 渐近可得
\[
H(\nu_{m,q})
=
m\log\varphi
-\frac12\log 5
+\bigl(1-\theta(q)\bigr)\log\varphi
+h_2(\theta(q))
+o(1).
\]
另一方面，由定理 \ref{thm:killo-fold-bin-escort-onebit-fisher}，
\[
\|\pi_{m,q}-\nu_{m,q}\|_{\mathrm{TV}}
=
O\!\Bigl(\Bigl(\frac{\varphi}{2}\Bigr)^m\Bigr).
\]
再用 Fannes--Audenaert 不等式与 \(|X_m|=F_{m+2}=e^{O(m)}\)，得到
\[
|H(\pi_{m,q})-H(\nu_{m,q})|
\le
\|\pi_{m,q}-\nu_{m,q}\|_{\mathrm{TV}}\log|X_m|
+h_2\!\bigl(\|\pi_{m,q}-\nu_{m,q}\|_{\mathrm{TV}}\bigr)
=o(1),
\]
故 Shannon 情形同样成立。

最后，上述两式的首项均为 \(m\log\varphi\)，故对一切固定 \(\alpha>0\) 都有
\[
\frac{1}{m}H_\alpha(\pi_{m,q})\to\log\varphi.
\]
\end{proof}

\endinput


\begin{theorem}[bin-fold 输出分布的全 Rényi 散度极限（对稳定层均匀）]\label{thm:fold-bin-renyi-divergence-limit}
令 \(\mu_m(w):=d^{\mathrm{bin}}_m(w)/2^m\) 为 bin-fold 的推前输出分布，\(u_m(w):=1/|X_m|\) 为稳定层均匀分布。
对任意 \(\alpha>0\) 定义 Rényi 散度 \(D_\alpha(\mu_m\Vert u_m)\)（\(\alpha=1\) 取 KL 散度的连续延拓）。
则极限
\[
D_\alpha^\infty:=\lim_{m\to\infty}D_\alpha(\mu_m\Vert u_m)
\]
存在，且给出闭式表达：
\begin{enumerate}
  \item 若 \(\alpha\neq 1\)，则
  \[
  D_\alpha^\infty=\frac{(2\alpha-2)\log\varphi+\log(\varphi+\varphi^{-\alpha})-\alpha\log\sqrt{5}}{\alpha-1}.
  \]
  \item 若 \(\alpha=1\)，则
  \[
  D_1^\infty=\frac{\varphi}{\sqrt5}\log\Bigl(\frac{\varphi^2}{\sqrt5}\Bigr)+\frac{1}{\varphi\sqrt5}\log\Bigl(\frac{\varphi}{\sqrt5}\Bigr).
  \]
\end{enumerate}
并且对每个固定 \(\alpha\)，收敛速度为
\[
D_\alpha(\mu_m\Vert u_m)=D_\alpha^\infty+O\bigl((\varphi/2)^m\bigr).
\]
\leanverified{paper\_fold\_bin\_renyi\_divergence\_limit}
\end{theorem}

\begin{proof}
当 \(\alpha\neq 1\) 时，由定义
\[
D_\alpha(\mu_m\Vert u_m)
=\frac{1}{\alpha-1}\log\sum_{w\in X_m}\mu_m(w)^\alpha\,u_m(w)^{1-\alpha}.
\]
代入 \(\mu_m(w)=d^{\mathrm{bin}}_m(w)/2^m\)、\(u_m(w)=1/|X_m|\) 得
\[
D_\alpha(\mu_m\Vert u_m)
=\frac{1}{\alpha-1}\log\Bigl(|X_m|^{\alpha-1}\cdot 2^{-\alpha m}\cdot S_\alpha^{\mathrm{bin}}(m)\Bigr).
\]
由 \(|X_m|=F_{m+2}\) 与命题 \ref{prop:fold-bin-power-sum-asymptotic}，
\[
2^{-\alpha m}S_\alpha^{\mathrm{bin}}(m)
=\varphi^{-\alpha m}\Bigl(F_{m+1}+\varphi^{-\alpha}F_m\Bigr)\Bigl(1+O\bigl((\varphi/2)^m\bigr)\Bigr).
\]
因此
\[
D_\alpha(\mu_m\Vert u_m)
=\frac{1}{\alpha-1}\log\Bigl(F_{m+2}^{\alpha-1}\varphi^{-\alpha m}\bigl(F_{m+1}+\varphi^{-\alpha}F_m\bigr)\Bigr)
O\bigl((\varphi/2)^m\bigr).
\]
将 Binet 渐近 \(F_n=\varphi^n/\sqrt5+O(\varphi^{-n})\) 代入并整理，可得极限常数为
\[
\frac{1}{\alpha-1}\Bigl((2\alpha-2)\log\varphi+\log(\varphi+\varphi^{-\alpha})-\alpha\log\sqrt5\Bigr),
\]
即所示 \(D_\alpha^\infty\)。误差界由命题 \ref{prop:fold-bin-power-sum-asymptotic} 的相对误差与 Binet 的指数余项合并得到。
\par
当 \(\alpha=1\) 时，\(D_1(\mu_m\Vert u_m)=D_{\mathrm{KL}}(\mu_m\Vert u_m)\)，其极限由命题 \ref{prop:fold-bin-kl-constant} 给出；同时由两态极限比值 \(\mu_m/u_m\Rightarrow\{\varphi^2/\sqrt5,\varphi/\sqrt5\}\) 亦可直接得到所示闭式。证毕。
\end{proof}

\begin{proposition}[bin-fold 的 \texorpdfstring{$f$}{f}-散度普适塌缩（Csiszár 指标）]\label{prop:fold-bin-f-divergence-collapse}
\leanverified{paper\_fold\_bin\_f\_divergence\_collapse}
令 \(f:(0,\infty)\to\RR\) 为凸函数且 \(f(1)=0\)，并定义 Csiszár \(f\)-散度
\[
D_f(\mu_m\Vert u_m):=\sum_{w\in X_m}u_m(w)\,f\Bigl(\frac{\mu_m(w)}{u_m(w)}\Bigr).
\]
则极限 \(D_f^\infty:=\lim_{m\to\infty}D_f(\mu_m\Vert u_m)\) 存在，并由两点闭式给出
\[
D_f^\infty=\frac{1}{\varphi}\,f\Bigl(\frac{\varphi^2}{\sqrt5}\Bigr)+\frac{1}{\varphi^2}\,f\Bigl(\frac{\varphi}{\sqrt5}\Bigr).
\]
并且对固定 \(f\)，有收敛速度 \(D_f(\mu_m\Vert u_m)=D_f^\infty+O\bigl((\varphi/2)^m\bigr)\)。
\end{proposition}

\begin{proof}
注意
\[
\frac{\mu_m(w)}{u_m(w)}=\frac{d^{\mathrm{bin}}_m(w)}{2^m}\cdot |X_m|
=d^{\mathrm{bin}}_m(w)\,\frac{F_{m+2}}{2^m}.
\]
由定理 \ref{thm:fold-bin-two-state-asymptotic} 与 Binet 渐近得一致展开
\[
\frac{\mu_m(w)}{u_m(w)}
=\frac{\varphi^{2-w_m}}{\sqrt5}+O\bigl((\varphi/2)^m\bigr),
\qquad (w\in X_m),
\]
并且在均匀测度 \(u_m\) 下，
\(
u_m(w_m=0)=F_{m+1}/F_{m+2}\to 1/\varphi
\)、
\(
u_m(w_m=1)=F_m/F_{m+2}\to 1/\varphi^2
\)。
于是
\[
D_f(\mu_m\Vert u_m)
=\frac{F_{m+1}}{F_{m+2}}\,f\Bigl(\frac{\varphi^2}{\sqrt5}+O\bigl((\varphi/2)^m\bigr)\Bigr)
\;+\;\frac{F_m}{F_{m+2}}\,f\Bigl(\frac{\varphi}{\sqrt5}+O\bigl((\varphi/2)^m\bigr)\Bigr).
\]
由凸函数在紧区间上的一致连续性即可令 \(m\to\infty\) 得到所示两点极限；误差阶同理由一致展开与 \(F_{m+1}/F_{m+2}-1/\varphi=O(\varphi^{-2m})\) 合并得到。证毕。
\end{proof}

\leanverified{paper\_fold\_bin\_phi\_identifiable\_power\_divergence}
\begin{proposition}[黄金参数的可辨识性：幂散度极限常数的严格单调反演]\label{prop:fold-bin-phi-identifiable-power-divergence}
对任意 \(\alpha>1\)，取严格凸函数
\[
f_\alpha(t):=t^\alpha-1-\alpha(t-1)\qquad(t>0),
\]
并记其 \(f\)-散度极限为 \(D_{f_\alpha}^\infty=\lim_{m\to\infty}D_{f_\alpha}(\mu_m\Vert u_m)\)（命题 \ref{prop:fold-bin-f-divergence-collapse}）。
则有闭式
\[
D_{f_\alpha}^\infty
=\frac{1}{5^{\alpha/2}}\Bigl(\varphi^{2\alpha-1}+\varphi^{\alpha-2}\Bigr)-1.
\]
将右端视为关于参数 \(\xi>1\) 的函数
\[
\mathcal{D}_\alpha(\xi):=\frac{1}{5^{\alpha/2}}\Bigl(\xi^{2\alpha-1}+\xi^{\alpha-2}\Bigr)-1,
\]
则 \(\mathcal{D}_\alpha:(1,\infty)\to\RR\) 严格递增；因此由单一标量 \(D_{f_\alpha}^\infty\) 可唯一反演 \(\xi\)，从而在本体系中可唯一恢复 \(\xi=\varphi\)。
\end{proposition}
\begin{proof}
由定义 \(\sum_{w\in X_m}u_m(w)\,\frac{\mu_m(w)}{u_m(w)}=\sum_w \mu_m(w)=1\)，故对任意 \(m\) 有
\[
D_{f_\alpha}(\mu_m\Vert u_m)=\sum_{w}u_m(w)\Bigl(\frac{\mu_m(w)}{u_m(w)}\Bigr)^\alpha-1.
\]
令 \(m\to\infty\)，由命题 \ref{prop:fold-bin-f-divergence-collapse} 的两点极限律可得
\[
D_{f_\alpha}^\infty
=\frac{1}{\varphi}\Bigl(\frac{\varphi^2}{\sqrt5}\Bigr)^\alpha+\frac{1}{\varphi^2}\Bigl(\frac{\varphi}{\sqrt5}\Bigr)^\alpha-1
=\frac{1}{5^{\alpha/2}}\Bigl(\varphi^{2\alpha-1}+\varphi^{\alpha-2}\Bigr)-1.
\]
对 \(\xi>1\)，\(\mathcal{D}_\alpha\) 的导数为
\[
\mathcal{D}_\alpha'(\xi)=\frac{1}{5^{\alpha/2}}\Bigl((2\alpha-1)\xi^{2\alpha-2}+(\alpha-2)\xi^{\alpha-3}\Bigr).
\]
当 \(\alpha\ge 2\) 时上式显然为正。
当 \(1<\alpha<2\) 时，因 \(\xi^{\alpha+1}\ge 1\) 且
\(
(2\alpha-1)-(2-\alpha)=3(\alpha-1)>0
\)，故
\[
(2\alpha-1)\xi^{2\alpha-2}+(\alpha-2)\xi^{\alpha-3}
\ge \xi^{\alpha-3}\bigl((2\alpha-1)\xi^{\alpha+1}-(2-\alpha)\bigr)
>0.
\]
因此 \(\mathcal{D}_\alpha'(\xi)>0\) 对一切 \(\xi>1\) 成立，\(\mathcal{D}_\alpha\) 严格递增，从而可逆。证毕。
\end{proof}

\begin{proposition}[规范体积密度的两尺度展开（bin-fold）]\label{prop:fold-bin-gauge-density-two-scale}
令 \(\mathcal{G}^{\mathrm{bin}}_m\) 为 bin-fold 的 fold-gauge 群体积
\(
|\mathcal{G}^{\mathrm{bin}}_m|=\prod_{w\in X_m} d^{\mathrm{bin}}_m(w)!
\)
并定义密度 \(g_m:=2^{-m}\log|\mathcal{G}^{\mathrm{bin}}_m|\)。
记
\(
\kappa_m:=\mathbb{E}_{\mu_m}\bigl[\log d^{\mathrm{bin}}_m(W_m)\bigr]
\)
（即命题 \ref{prop:fold-bin-kl-constant} 中的纤维信息项）。
则当 \(m\to\infty\) 时有二尺度展开
\[
g_m=\kappa_m-1+\frac{(\varphi/2)^m}{2\sqrt5}\Bigl(\varphi^2 m\log\!\Bigl(\frac{2}{\varphi}\Bigr)+\varphi^2\log(2\pi)-\log\varphi\Bigr)+O\bigl((\varphi/2)^m\bigr).
\]
\leanverified{paper\_fold\_bin\_gauge\_density\_two\_scale}
\end{proposition}

\begin{proof}
由定理 \ref{thm:fold-bin-gauge-volume-stirling-second-order}，
\[
g_m=\kappa_m-1+\frac{1}{2^{m+1}}\Bigl(|X_m|\log(2\pi)+\sum_{w\in X_m}\log d^{\mathrm{bin}}_m(w)\Bigr)+R_m,
\qquad 0\le R_m\le \frac{|X_m|}{12\cdot 2^m}.
\]
由定理 \ref{thm:fold-bin-two-state-asymptotic} 的一致对数展开，
\[
\log d^{\mathrm{bin}}_m(w)=m\log\!\Bigl(\frac{2}{\varphi}\Bigr)-w_m\log\varphi+O\bigl((\varphi/2)^m\bigr),
\qquad (w\in X_m),
\]
求和并用
\(
\card\{w\in X_m:w_m=0\}=F_{m+1},\ \card\{w\in X_m:w_m=1\}=F_m
\)
得
\[
\sum_{w\in X_m}\log d^{\mathrm{bin}}_m(w)
=m\log\!\Bigl(\frac{2}{\varphi}\Bigr)|X_m|-\log\varphi\cdot F_m+O\bigl(|X_m|(\varphi/2)^m\bigr).
\]
代回并用 Binet 渐近 \(F_{m+2}=|X_m|=\varphi^{m+2}/\sqrt5+O(\varphi^{-m})\)、\(F_m=\varphi^m/\sqrt5+O(\varphi^{-m})\) 得
\[
\frac{|X_m|}{2^{m+1}}
=\frac{\varphi^2}{2\sqrt5}\Bigl(\frac{\varphi}{2}\Bigr)^m+O\bigl((\varphi/2)^m\bigr),
\qquad
\frac{F_m}{2^{m+1}}
=\frac{1}{2\sqrt5}\Bigl(\frac{\varphi}{2}\Bigr)^m+O\bigl((\varphi/2)^m\bigr),
\]
并由 \(R_m=O(|X_m|/2^m)=O((\varphi/2)^m)\) 合并得到所示展开。证毕。
\end{proof}

\begin{proposition}[bin-fold 退化谱的尺度极限：尾计数与连续容量曲线的两折点]\label{prop:fold-bin-degeneracy-tail-capacity-kinks}
对每个 \(m\ge 1\)，定义 bin-fold 退化度尾计数函数
\[
N^{\mathrm{bin}}_m(t):=\bigl|\{w\in X_m:\ d^{\mathrm{bin}}_m(w)\ge t\}\bigr|,\qquad t\ge 0,
\]
以及连续化容量曲线
\[
\mathcal{C}^{\mathrm{bin,cont}}_m(T):=\sum_{w\in X_m}\min\bigl(d^{\mathrm{bin}}_m(w),T\bigr),\qquad T\ge 0.
\]
则对所有 \(T\ge 0\) 有严格恒等式
\[
\mathcal{C}^{\mathrm{bin,cont}}_m(T)=\int_{0}^{T}N^{\mathrm{bin}}_m(t)\,\dd t.
\]
对固定 \(s>0\) 定义尺度化过程
\[
\widehat N^{\mathrm{bin}}_m(s):=\frac{1}{|X_m|}\,N^{\mathrm{bin}}_m\!\Bigl(s\Bigl(\frac{2}{\varphi}\Bigr)^m\Bigr),
\qquad
\widehat C^{\mathrm{bin}}_m(s):=2^{-m}\,\mathcal{C}^{\mathrm{bin,cont}}_m\!\Bigl(s\Bigl(\frac{2}{\varphi}\Bigr)^m\Bigr).
\]
则对每个 \(s\notin\{\varphi^{-1},1\}\) 有点态极限
\[
\widehat N^{\mathrm{bin}}_m(s)\longrightarrow
\begin{cases}
1, & 0<s<\varphi^{-1},\\[2pt]
\varphi^{-1}, & \varphi^{-1}<s<1,\\[2pt]
0, & s>1,
\end{cases}
\]
并且 \(\widehat C^{\mathrm{bin}}_m\) 在每个紧集上一致收敛到分段线性函数
\[
\widehat C^{\mathrm{bin}}_\infty(s)
=\frac{\varphi}{\sqrt5}\min\{1,s\}+\frac{1}{\sqrt5}\min\{\varphi^{-1},s\}.
\]
因此在 bin-fold 极限下，尺度化“可见容量—尾谱”出现两处刚性折点 \(s=\varphi^{-1}\) 与 \(s=1\)。
\leanverified{paper\_fold\_bin\_degeneracy\_tail\_capacity\_kinks}
\end{proposition}
\begin{proof}
容量—尾函数恒等式由 \(\min(d,T)=\int_0^T\ind_{\{t\le d\}}\,\dd t\) 对 \(w\) 求和并交换求和与积分直接得到。
\par
令 \(t_m(s):=s(2/\varphi)^m\)。
由定理 \ref{thm:fold-bin-two-state-asymptotic} 的一致渐近，
\[
\Bigl(\frac{\varphi}{2}\Bigr)^m d^{\mathrm{bin}}_m(w)=\varphi^{-w_m}+O\bigl((\varphi/2)^m\bigr),
\qquad (w\in X_m),
\]
故当 \(s\notin\{\varphi^{-1},1\}\) 固定时，充分大 \(m\) 上不等式
\(
d^{\mathrm{bin}}_m(w)\ge t_m(s)
\)
与末位条件等价：
\begin{itemize}
  \item 若 \(0<s<\varphi^{-1}\)，则对一切 \(w\in X_m\) 成立；
  \item 若 \(\varphi^{-1}<s<1\)，则当且仅当 \(w_m=0\)；
  \item 若 \(s>1\)，则对一切 \(w\in X_m\) 不成立。
\end{itemize}
因此
\[
N^{\mathrm{bin}}_m(t_m(s))
=
\begin{cases}
|X_m|,& 0<s<\varphi^{-1},\\
|A_0|,& \varphi^{-1}<s<1,\\
0,& s>1,
\end{cases}
\qquad A_0=\{w:\ w_m=0\},
\]
从而
\(
\widehat N^{\mathrm{bin}}_m(s)\to 1
\)
或
\(
\widehat N^{\mathrm{bin}}_m(s)\to |A_0|/|X_m|=F_{m+1}/F_{m+2}\to \varphi^{-1}
\)，得到所示点态极限。
\par
对容量曲线，把 \(X_m\) 按 \(w_m\in\{0,1\}\) 分块并用两态一致渐近得到
\[
\mathcal{C}^{\mathrm{bin,cont}}_m(t_m(s))
=\sum_{w_m=0}\min\!\Bigl(\Bigl(\frac{2}{\varphi}\Bigr)^m+O(1),\ s\Bigl(\frac{2}{\varphi}\Bigr)^m\Bigr)
\;+\;
\sum_{w_m=1}\min\!\Bigl(\varphi^{-1}\Bigl(\frac{2}{\varphi}\Bigr)^m+O(1),\ s\Bigl(\frac{2}{\varphi}\Bigr)^m\Bigr).
\]
对固定 \(s\) 与充分大 \(m\)，两处 \(\min\) 的主阶分别为 \((2/\varphi)^m\min\{1,s\}\) 与 \((2/\varphi)^m\min\{\varphi^{-1},s\}\)，并且误差为 \(O(|X_m|)\)。
因此
\[
\widehat C^{\mathrm{bin}}_m(s)
=2^{-m}\Bigl(\frac{2}{\varphi}\Bigr)^m\Bigl(|A_0|\min\{1,s\}+|A_1|\min\{\varphi^{-1},s\}\Bigr)+O\!\Bigl(\frac{|X_m|}{2^m}\Bigr),
\]
其中 \(|A_0|=F_{m+1}\)、\(|A_1|=F_m\)。
由 Binet 渐近 \(F_{m+1}/\varphi^m\to \varphi/\sqrt5\)、\(F_m/\varphi^m\to 1/\sqrt5\)，并注意
\(
2^{-m}(2/\varphi)^m=\varphi^{-m}
\)，得到
\[
\widehat C^{\mathrm{bin}}_m(s)\to \frac{\varphi}{\sqrt5}\min\{1,s\}+\frac{1}{\sqrt5}\min\{\varphi^{-1},s\}.
\]
误差项 \(O(|X_m|/2^m)=O((\varphi/2)^m)\) 在每个紧集上一致，从而得到一致收敛。证毕。
\end{proof}

\begin{theorem}[两折点尾谱的 Mellin 矩谱刻画与唯一性]\label{thm:killo-fold-bin-two-kink-mellin-uniqueness}
\leanverified{paper\_killo\_fold\_bin\_two\_kink\_mellin\_uniqueness}
沿用命题 \ref{prop:fold-bin-degeneracy-tail-capacity-kinks} 的记号，并记
\[
\widehat N_m:=\widehat N_m^{\mathrm{bin}},
\qquad
\widehat C_m:=\widehat C_m^{\mathrm{bin}}.
\]
则对任意 \(q>0\)，归一化矩谱极限存在且满足
\[
\lim_{m\to\infty}
\frac{S_q^{\mathrm{bin}}(m)}{(2/\varphi)^{qm}|X_m|}
=
q\int_{0}^{\infty}s^{q-1}\widehat N_\infty(s)\,\dd s
=
\varphi^{-1}+\varphi^{-q-2},
\]
其中
\[
\widehat N_\infty(s)=
\begin{cases}
1, & 0\le s<\varphi^{-1},\\[2pt]
\varphi^{-1}, & \varphi^{-1}\le s<1,\\[2pt]
0, & s\ge 1.
\end{cases}
\]
反之，若 \(N:[0,\infty)\to[0,1]\) 为右连续、单调不增的函数，满足
\[
N(0)=1,
\qquad
N(s)=0\quad(s>1),
\]
且对一切 \(q>0\) 都有
\[
q\int_{0}^{\infty}s^{q-1}N(s)\,\dd s
=
\varphi^{-1}+\varphi^{-q-2},
\]
则必有 \(N=\widehat N_\infty\)。
\end{theorem}
\begin{proof}
由命题 \ref{prop:fold-moment-mellin-stieltjes}，
\[
S_q^{\mathrm{bin}}(m)=q\int_{0}^{\infty}t^{q-1}N_m^{\mathrm{bin}}(t)\,\dd t.
\]
作变量代换 \(t=s(2/\varphi)^m\)，得到
\[
\frac{S_q^{\mathrm{bin}}(m)}{(2/\varphi)^{qm}|X_m|}
=
q\int_{0}^{\infty}s^{q-1}\widehat N_m(s)\,\dd s.
\]
另一方面，由定理 \ref{thm:fold-bin-two-state-asymptotic} 的一致渐近，
\[
\max_{w\in X_m}\Bigl(\frac{\varphi}{2}\Bigr)^m d_m^{\mathrm{bin}}(w)=1+o(1),
\]
故对所有充分大的 \(m\) 有 \(\widehat N_m(s)=0\) 当 \(s>2\)。
再结合 \(0\le \widehat N_m(s)\le 1\) 与命题 \ref{prop:fold-bin-degeneracy-tail-capacity-kinks} 的点态极限，可由支配收敛定理得到
\[
\lim_{m\to\infty}
\frac{S_q^{\mathrm{bin}}(m)}{(2/\varphi)^{qm}|X_m|}
=
q\int_{0}^{\infty}s^{q-1}\widehat N_\infty(s)\,\dd s.
\]
将 \(\widehat N_\infty\) 的显式表达代入可得
\[
q\int_{0}^{\varphi^{-1}}s^{q-1}\,\dd s
+
q\varphi^{-1}\int_{\varphi^{-1}}^{1}s^{q-1}\,\dd s
=
\varphi^{-q}+\varphi^{-1}\bigl(1-\varphi^{-q}\bigr)
=
\varphi^{-1}+\varphi^{-q-2},
\]
即得第一部分。

再证唯一性。
由右连续与单调性，可定义有限 Borel 测度 \(\nu\) 于 \([0,1]\) 上，使
\[
N(s)=\nu([s,1])\qquad (0\le s\le 1).
\]
对该尾函数作分部积分，得到
\[
q\int_{0}^{1}s^{q-1}N(s)\,\dd s=\int_{0}^{1}x^q\,\dd \nu(x).
\]
因此假设等价于
\[
\int_{0}^{1}x^q\,\dd \nu(x)=\varphi^{-1}+\varphi^{-q-2}
\qquad (\forall q>0).
\]
记
\[
a:=\varphi^{-1},
\qquad
b:=\varphi^{-2},
\qquad
c:=\varphi^{-1}.
\]
令 \(q\to\infty\)，由 \(x^q\to \mathbf 1_{\{x=1\}}\) 且 \(0\le x^q\le 1\)，支配收敛给出
\[
\nu(\{1\})=\lim_{q\to\infty}\int_{0}^{1}x^q\,\dd \nu(x)=a.
\]
于是可写
\[
\nu=a\delta_1+\eta,
\]
其中 \(\eta\) 为总质量 \(b\) 的正测度，并满足
\[
\int_{0}^{1}x^q\,\dd \eta(x)=bc^q
\qquad (\forall q>0).
\]
若 \(\eta((c,1])>0\)，则存在 \(\varepsilon>0\) 使 \(\eta([c+\varepsilon,1])>0\)，从而
\[
bc^q=\int x^q\,\dd \eta(x)\ge (c+\varepsilon)^q\eta([c+\varepsilon,1]),
\]
这在 \(q\to\infty\) 时矛盾。
故 \(\eta\) 的支撑包含于 \([0,c]\)。

令 \(\widetilde \eta:=b^{-1}\eta\)。
则 \(\widetilde \eta\) 为 \([0,c]\) 上的概率测度，且
\[
\int_{0}^{c}\Bigl(\frac{x}{c}\Bigr)^q\,\dd \widetilde \eta(x)=1
\qquad (\forall q>0).
\]
令 \(q\to\infty\)，再次用支配收敛得
\[
\widetilde \eta(\{c\})=1,
\]
故 \(\widetilde \eta=\delta_c\)，即 \(\eta=b\delta_c\)。
因此
\[
\nu=\varphi^{-1}\delta_1+\varphi^{-2}\delta_{\varphi^{-1}}.
\]
其尾函数正是 \(\widehat N_\infty\)，从而 \(N=\widehat N_\infty\)。
\end{proof}

\begin{theorem}[连续容量反函数的二段精确律]\label{thm:killo-fold-bin-capacity-inverse-two-regime}
\leanverified{paper\_killo\_fold\_bin\_capacity\_inverse\_two\_regime}
继续记
\[
\widehat C_m:=\widehat C_m^{\mathrm{bin}}.
\]
对任意 \(c\in(0,1)\)，定义缩放阈值
\[
\mathcal B_m(c):=\inf\{s\ge 0:\ \widehat C_m(s)\ge c\}.
\]
则 \(\mathcal B_m(c)\to \mathcal B_\infty(c)\)，其中
\[
\mathcal B_\infty(c)=\widehat C_\infty^{-1}(c)
=
\begin{cases}
\dfrac{\sqrt5}{\varphi^2}\,c,
& 0<c\le \dfrac{\varphi}{\sqrt5},\\[2ex]
\dfrac{\sqrt5}{\varphi}\,c-\varphi^{-2},
& \dfrac{\varphi}{\sqrt5}\le c<1.
\end{cases}
\]
若再定义未缩放连续容量阈值
\[
T_m(c):=\inf\Bigl\{T\ge 0:\ 2^{-m}\mathcal C_m^{\mathrm{bin,cont}}(T)\ge c\Bigr\},
\]
则有精确关系
\[
T_m(c)=\Bigl(\frac{2}{\varphi}\Bigr)^m\mathcal B_m(c),
\]
从而
\[
T_m(c)=\Bigl(\frac{2}{\varphi}\Bigr)^m\bigl(\mathcal B_\infty(c)+o(1)\bigr).
\]
相应地，若记连续容量的对数预算为
\[
B_m(c):=\log_2 T_m(c),
\]
则
\[
B_m(c)=m\log_2\!\Bigl(\frac{2}{\varphi}\Bigr)+\log_2\mathcal B_\infty(c)+o(1).
\]
\end{theorem}
\begin{proof}
由命题 \ref{prop:fold-bin-degeneracy-tail-capacity-kinks}，
\[
\widehat C_m(s)\to \widehat C_\infty(s)
\]
在每个紧区间上一致，并且
\[
\widehat C_\infty(s)=\frac{\varphi}{\sqrt5}\min\{1,s\}+\frac{1}{\sqrt5}\min\{\varphi^{-1},s\}.
\]
将其写成分段式即
\[
\widehat C_\infty(s)=
\begin{cases}
\dfrac{\varphi^2}{\sqrt5}\,s, & 0\le s\le \varphi^{-1},\\[2ex]
\dfrac{\varphi}{\sqrt5}\,s+\dfrac{1}{\varphi\sqrt5}, & \varphi^{-1}\le s\le 1,\\[2ex]
1, & s\ge 1.
\end{cases}
\]
故 \(\widehat C_\infty\) 在 \([0,1]\) 上连续且严格递增。

固定 \(c\in(0,1)\)，令 \(s_\ast:=\widehat C_\infty^{-1}(c)\in(0,1)\)。
任取 \(\varepsilon>0\) 充分小，使 \(0<s_\ast-\varepsilon<s_\ast+\varepsilon<1\)。
由严格单调性，
\[
\widehat C_\infty(s_\ast-\varepsilon)<c<\widehat C_\infty(s_\ast+\varepsilon).
\]
再由 \(\widehat C_m\to \widehat C_\infty\) 在 \([0,1]\) 上一致，得到对所有充分大的 \(m\)，
\[
\widehat C_m(s_\ast-\varepsilon)<c<\widehat C_m(s_\ast+\varepsilon).
\]
按 \(\mathcal B_m(c)\) 的定义，必有
\[
s_\ast-\varepsilon\le \mathcal B_m(c)\le s_\ast+\varepsilon.
\]
由于 \(\varepsilon\) 任意，故 \(\mathcal B_m(c)\to s_\ast=\mathcal B_\infty(c)\)。
对上述分段线性函数逐段求逆，即得 \(\mathcal B_\infty(c)\) 的显式表达。

再由
\[
\widehat C_m(s)=2^{-m}\mathcal C_m^{\mathrm{bin,cont}}\!\Bigl(s\Bigl(\frac{2}{\varphi}\Bigr)^m\Bigr)
\]
可知 \(T=s(2/\varphi)^m\) 给出阈值变量的精确缩放，因此
\[
T_m(c)=\Bigl(\frac{2}{\varphi}\Bigr)^m\mathcal B_m(c).
\]
将 \(\mathcal B_m(c)\to \mathcal B_\infty(c)\) 代入即得
\[
T_m(c)=\Bigl(\frac{2}{\varphi}\Bigr)^m\bigl(\mathcal B_\infty(c)+o(1)\bigr).
\]
最后取二进制对数便得到
\[
B_m(c)=m\log_2\!\Bigl(\frac{2}{\varphi}\Bigr)+\log_2\mathcal B_\infty(c)+o(1).
\]
\end{proof}

\endinput


\begin{corollary}[由 fold-gauge 体积与纤维信息项内生恢复 \texorpdfstring{\(2\pi\)}{2pi}]\label{cor:fold-bin-recover-2pi}
\leanverified{paper\_fold\_bin\_recover\_2pi}
沿用命题 \ref{prop:fold-bin-gauge-density-two-scale} 的 \(g_m,\kappa_m\) 记号。
定义由折叠数据直接给出的常数估计序列
\[
C_m
:=\exp\!\Biggl(
\frac{2^{m+1}}{|X_m|}\bigl(g_m-\kappa_m+1\bigr)
-\frac{1}{|X_m|}\sum_{w\in X_m}\log d^{\mathrm{bin}}_m(w)
\Biggr).
\]
则极限存在并满足
\[
\lim_{m\to\infty}C_m=2\pi.
\]
并且收敛具有指数速度：存在常数 \(c>0\) 使对所有充分大的 \(m\) 有
\[
\bigl|\log C_m-\log(2\pi)\bigr|\le c\Bigl(\frac{\varphi}{2}\Bigr)^m.
\]
\end{corollary}
\begin{proof}
对任意整数 \(d\ge 1\)，带余项的 Stirling 公式给出
\[
\log(d!)=d\log d-d+\frac12\log(2\pi d)+\varepsilon(d),
\qquad |\varepsilon(d)|\le \frac{1}{12d}.
\]
对 \(d=d^{\mathrm{bin}}_m(w)\) 求和并除以 \(2^m\)，并用 \(\sum_w d^{\mathrm{bin}}_m(w)=2^m\)，得到
\[
g_m
=\kappa_m-1+\frac{1}{2^{m+1}}\sum_{w\in X_m}\log(2\pi d^{\mathrm{bin}}_m(w))
\;+\;\frac{1}{2^m}\sum_{w\in X_m}\varepsilon\bigl(d^{\mathrm{bin}}_m(w)\bigr).
\]
将
\(
\sum_w \log(2\pi d)=|X_m|\log(2\pi)+\sum_w \log d
\)
代回并整理，即得
\[
\log C_m=\log(2\pi)+\frac{2}{|X_m|}\sum_{w\in X_m}\varepsilon\bigl(d^{\mathrm{bin}}_m(w)\bigr).
\]
由定理 \ref{thm:fold-bin-two-state-asymptotic} 的一致渐近，存在 \(m_0\) 与常数 \(c_0>0\) 使当 \(m\ge m_0\) 时对所有 \(w\in X_m\) 有
\(
d^{\mathrm{bin}}_m(w)\ge c_0(2/\varphi)^m
\)。
因此
\[
\Bigl|\frac{2}{|X_m|}\sum_{w}\varepsilon\bigl(d^{\mathrm{bin}}_m(w)\bigr)\Bigr|
\le \frac{2}{|X_m|}\sum_{w}\frac{1}{12\,d^{\mathrm{bin}}_m(w)}
\le \frac{1}{6}\cdot \frac{1}{\min_w d^{\mathrm{bin}}_m(w)}
\le c\Bigl(\frac{\varphi}{2}\Bigr)^m,
\]
从而 \(\log C_m\to\log(2\pi)\) 且给出指数收敛界。证毕。
\end{proof}

\begin{theorem}[折叠规范群的 Stirling--Bernoulli 层级：由折叠数据内生读出全部 Bernoulli 系数]\label{thm:fold-bin-gauge-constant-stirling-bernoulli-hierarchy}
沿用推论 \ref{cor:fold-bin-recover-2pi} 的 \(C_m\) 定义，并记 \(\log C_m\) 为其自然对数。
则存在完整渐近展开：对任意固定整数 \(R\ge 1\)，当 \(m\to\infty\) 时
\[
\boxed{\;
\log C_m
=\log(2\pi)+
\sum_{r=1}^{R}
\frac{B_{2r}}{r(2r-1)}
\Bigl(\varphi^{-1}+\varphi^{2r-3}\Bigr)
\Bigl(\frac{\varphi}{2}\Bigr)^{(2r-1)m}
\;+\;
O\!\Bigl(\Bigl(\frac{\varphi}{2}\Bigr)^{(2R+1)m}\Bigr)\; } ,
\]
其中 \(B_{2r}\) 为 Bernoulli 数。
因此对每个 \(r\ge 1\)，Bernoulli 系数 \(B_{2r}\) 可由折叠数据序列 \(\{\log C_m\}_{m\ge 1}\) 通过逐阶消元在极限中唯一恢复；
并且结合标准恒等式
\[
B_{2r}=(-1)^{r-1}\frac{2(2r)!}{(2\pi)^{2r}}\zeta(2r),
\]
可将全部 \(\zeta(2r)\) 视作该折叠规范常数的渐近谱信息。
\end{theorem}
\begin{proof}
采用 Stirling--Bernoulli 渐近展开（作为外部工具）：对任意固定 \(R\ge 1\) 与整数 \(d\ge 1\) 有
\[
\log(d!)=d\log d-d+\frac12\log(2\pi d)+
\sum_{r=1}^{R}\frac{B_{2r}}{2r(2r-1)}\,d^{-(2r-1)}
\;+\;
O\!\bigl(d^{-(2R+1)}\bigr).
\]
将其应用到 \(d=d^{\mathrm{bin}}_m(w)\) 并对 \(w\in X_m\) 求和，再除以 \(2^m\)，并按
\(
g_m:=2^{-m}\sum_w\log(d^{\mathrm{bin}}_m(w)!)
\)
与
\(
\kappa_m:=2^{-m}\sum_w d^{\mathrm{bin}}_m(w)\log d^{\mathrm{bin}}_m(w)
\)
整理，可得
\[
g_m
=\kappa_m-1+\frac{1}{2^{m+1}}\sum_{w\in X_m}\log(2\pi d^{\mathrm{bin}}_m(w))
\;+\;
\sum_{r=1}^{R}\frac{B_{2r}}{2r(2r-1)}\cdot
\frac{1}{2^{m}}\sum_{w\in X_m} d^{\mathrm{bin}}_m(w)^{-(2r-1)}
\;+\;
O\!\left(\frac{1}{2^{m}}\sum_{w\in X_m} d^{\mathrm{bin}}_m(w)^{-(2R+1)}\right).
\]
将
\(
\sum_w\log(2\pi d)=|X_m|\log(2\pi)+\sum_w\log d
\)
代回，并代入 \(\log C_m\) 的定义（推论 \ref{cor:fold-bin-recover-2pi}），主项严格抵消，仅留下 Bernoulli 级数的均匀平均：
\[
\log C_m
=\log(2\pi)+
\sum_{r=1}^{R}\frac{B_{2r}}{r(2r-1)}\cdot
\frac{1}{|X_m|}\sum_{w\in X_m} d^{\mathrm{bin}}_m(w)^{-(2r-1)}
\;+\;
O\!\left(\frac{1}{|X_m|}\sum_{w\in X_m} d^{\mathrm{bin}}_m(w)^{-(2R+1)}\right).
\]
\par
由定理 \ref{thm:fold-bin-two-state-asymptotic} 的一致渐近
\(
d^{\mathrm{bin}}_m(w)=\varphi^{-w_m}(2/\varphi)^m+O(1)
\)
以及末位分块 \(\#\{w\in X_m:w_m=0\}=F_{m+1}\)、\(\#\{w\in X_m:w_m=1\}=F_m\) 与 \(F_{m+1}/F_{m+2}\to\varphi^{-1}\)、\(F_m/F_{m+2}\to\varphi^{-2}\)，可得对每个固定 \(r\ge 1\) 有
\[
\frac{1}{|X_m|}\sum_{w\in X_m} d^{\mathrm{bin}}_m(w)^{-(2r-1)}
=\Bigl(\varphi^{-1}+\varphi^{2r-3}\Bigr)\Bigl(\frac{\varphi}{2}\Bigr)^{(2r-1)m}
\;+\;
O\!\Bigl(\Bigl(\frac{\varphi}{2}\Bigr)^{(2r+1)m}\Bigr).
\]
同理，
\(
|X_m|^{-1}\sum_w d^{-(2R+1)}
=O((\varphi/2)^{(2R+1)m})
\)。
代回即得所示展开与余项阶。最后两句关于“逐阶消元恢复 \(B_{2r}\)”与 \(\zeta(2r)\) 的编码为标准推论：由于各阶指数 \((\varphi/2)^{(2r-1)m}\) 相互分离，可通过依次去除低阶项并令 \(m\to\infty\) 唯一读出每个系数。证毕。
\end{proof}

\begin{proposition}[Bernoulli 系数的显式抽取算子：指数分离给出无外参的极限读出]\label{prop:fold-bin-gauge-bernoulli-extraction-operator}
\leanverified{paper\_fold\_bin\_gauge\_bernoulli\_extraction\_operator}
沿用定理 \ref{thm:fold-bin-gauge-constant-stirling-bernoulli-hierarchy} 的记号，并令
\[
q:=\frac{\varphi}{2}\in(0,1),\qquad
L_m:=\log C_m-\lim_{k\to\infty}\log C_k=\log C_m-\log(2\pi).
\]
递归定义系数读出：
\[
a_1:=\lim_{m\to\infty}q^{-m}L_m,
\]
以及对每个 \(r\ge 2\),
\[
a_r:=\lim_{m\to\infty} q^{-(2r-1)m}\left(L_m-\sum_{j=1}^{r-1} a_j\,q^{(2j-1)m}\right).
\]
则每个极限 \(a_r\) 均存在，并满足
\[
a_r=\frac{B_{2r}}{r(2r-1)}\Bigl(\varphi^{-1}+\varphi^{2r-3}\Bigr).
\]
从而全部 \(B_{2r}\)（并进而全部 \(\zeta(2r)\)）可由折叠输出序列 \(\{\log C_m\}_{m\ge 1}\) 逐阶确定性恢复；该恢复仅使用 \(C_m\) 序列本身及其极限 \(\lim_{m\to\infty}\log C_m\)，不依赖外部拟合参数。
\end{proposition}
\begin{proof}
由定理 \ref{thm:fold-bin-gauge-constant-stirling-bernoulli-hierarchy} 的展开，对任意固定 \(R\ge 1\) 有
\[
L_m=\sum_{r=1}^{R}\alpha_r\,q^{(2r-1)m}+O\!\bigl(q^{(2R+1)m}\bigr),
\qquad m\to\infty,
\]
其中
\[
\alpha_r:=\frac{B_{2r}}{r(2r-1)}\Bigl(\varphi^{-1}+\varphi^{2r-3}\Bigr).
\]
取 \(R=1\) 得
\(
q^{-m}L_m=\alpha_1+O(q^{2m})
\)，故极限 \(a_1=\lim_{m\to\infty}q^{-m}L_m\) 存在且 \(a_1=\alpha_1\)。
\par
设对某个 \(r\ge 2\) 已有 \(a_j=\alpha_j\)（\(1\le j\le r-1\)）。对定理展开取 \(R=r\)，移项并代入归纳假设，得
\[
L_m-\sum_{j=1}^{r-1}a_j\,q^{(2j-1)m}
=\alpha_r\,q^{(2r-1)m}+O\!\bigl(q^{(2r+1)m}\bigr).
\]
两边乘以 \(q^{-(2r-1)m}\) 得
\[
q^{-(2r-1)m}\left(L_m-\sum_{j=1}^{r-1}a_j\,q^{(2j-1)m}\right)
=\alpha_r+O(q^{2m}),
\]
从而极限 \(a_r\) 存在且 \(a_r=\alpha_r\)。最后由
\(
B_{2r}=(-1)^{r-1}\frac{2(2r)!}{(2\pi)^{2r}}\zeta(2r)
\)
即可逐阶恢复全部偶值 \(\zeta(2r)\)。证毕。
\end{proof}

\endinput

\begin{theorem}[Stirling--Bernoulli 指数层级的非 \(C\)-有限刚性]\label{thm:fold-bin-gauge-constant-not-c-finite}
\leanverified{paper\_fold\_bin\_gauge\_constant\_not\_c\_finite}
令
\[
u_m:=\log C_m-\log(2\pi),
\qquad
U(z):=\sum_{m\ge 1}u_m z^m.
\]
则序列 \((u_m)_{m\ge 1}\) 不满足任何有限阶常系数线性递推；等价地，普通生成函数 \(U(z)\) 不是有理函数。
\end{theorem}
\begin{proof}
设
\[
q:=\frac{\varphi}{2}\in(0,1),
\qquad
A_r:=\frac{B_{2r}}{r(2r-1)}\bigl(\varphi^{-1}+\varphi^{2r-3}\bigr),
\qquad
\beta_r:=q^{2r-1}.
\]
由定理 \ref{thm:fold-bin-gauge-constant-stirling-bernoulli-hierarchy}，对任意固定 \(R\ge 1\) 都有
\[
u_m=\sum_{r=1}^{R}A_r\beta_r^m+O(\beta_{R+1}^m),
\qquad m\to\infty.
\]
并且对每个 \(r\ge 1\)，Bernoulli 数 \(B_{2r}\neq 0\)，故 \(A_r\neq 0\)；又 \(\beta_r\) 随 \(r\) 严格递减，因此 \((u_m)\) 具有无穷多个彼此不同的指数尺度。

反设 \((u_m)\) 为 \(C\)-有限序列。
则存在有限个互异复数 \(\alpha_1,\dots,\alpha_J\neq 0\) 及多项式 \(P_1,\dots,P_J\)，使得对充分大的 \(m\) 有
\[
u_m=\sum_{j=1}^{J}P_j(m)\alpha_j^m.
\]
因此其指数尺度只能来自有限集合 \(\{|\alpha_1|,\dots,|\alpha_J|\}\)。
这与上式给出的无限严格递减尺度族 \(\{\beta_r\}_{r\ge 1}\) 矛盾。
故 \((u_m)\) 不可能是 \(C\)-有限序列。

而常数系数线性递推序列当且仅当其普通生成函数为有理函数，因此 \(U(z)\) 非有理。证毕。
\end{proof}

\begin{theorem}[Bernoulli 系数与 \texorpdfstring{$\zeta(2r)$}{zeta(2r)} 的显式极限恢复]\label{thm:fold-bin-gauge-constant-zeta-even-recovery-limit}
\leanverified{paper\_fold\_bin\_gauge\_constant\_zeta\_even\_recovery\_limit}
令
\[
q:=\frac{\varphi}{2}\in(0,1),
\qquad
u_m:=\log C_m-\log(2\pi),
\]
并对每个整数 \(r\ge 1\) 定义
\[
A_r:=\frac{B_{2r}}{r(2r-1)}\bigl(\varphi^{-1}+\varphi^{2r-3}\bigr).
\]
则成立显式极限公式
\[
\boxed{\;
A_r=
\lim_{m\to\infty}
q^{-(2r-1)m}
\left(
u_m-\sum_{j=1}^{r-1}A_j\,q^{(2j-1)m}
\right).\;}
\]
因此
\[
\boxed{\;
B_{2r}
=
\frac{r(2r-1)}{\varphi^{-1}+\varphi^{2r-3}}
\lim_{m\to\infty}
q^{-(2r-1)m}
\left(
u_m-\sum_{j=1}^{r-1}A_j\,q^{(2j-1)m}
\right).\;}
\]
再结合恒等式
\[
B_{2r}=(-1)^{r-1}\frac{2(2r)!}{(2\pi)^{2r}}\zeta(2r),
\]
得到
\[
\boxed{\;
\zeta(2r)
=
(-1)^{r-1}\frac{(2\pi)^{2r}}{2(2r)!}\,
\frac{r(2r-1)}{\varphi^{-1}+\varphi^{2r-3}}
\lim_{m\to\infty}
q^{-(2r-1)m}
\left(
u_m-\sum_{j=1}^{r-1}A_j\,q^{(2j-1)m}
\right).\;}
\]
\end{theorem}
\begin{proof}
由定理 \ref{thm:fold-bin-gauge-constant-stirling-bernoulli-hierarchy}，对任意固定 \(r\ge 1\) 有
\[
u_m=\sum_{j=1}^{r-1}A_j\,q^{(2j-1)m}+A_r\,q^{(2r-1)m}+O(q^{(2r+1)m}).
\]
移去前 \(r-1\) 层后乘以 \(q^{-(2r-1)m}\)，得到
\[
q^{-(2r-1)m}
\left(
u_m-\sum_{j=1}^{r-1}A_j\,q^{(2j-1)m}
\right)
=A_r+O(q^{2m}).
\]
令 \(m\to\infty\) 即得第一式。
第二式只需把 \(A_r\) 的定义整理为 \(B_{2r}\)；
第三式则由 Bernoulli 与 \(\zeta(2r)\) 的标准恒等式直接代入。
这也把命题 \ref{prop:fold-bin-gauge-bernoulli-extraction-operator} 的递归抽取，改写为显式的残差缩放极限。
证毕。
\end{proof}

\begin{proposition}[归一化规范亏损的 Stirling--Bernoulli 展开]\label{prop:killo-fold-bin-normalized-gauge-deficiency}
\leanverified{paper\_killo\_fold\_bin\_normalized\_gauge\_deficiency}
沿用命题 \ref{prop:fold-bin-gauge-density-two-scale} 与推论 \ref{cor:fold-bin-recover-2pi} 的记号，定义
\[
\Delta_m:=
\frac{2^{m+1}}{|X_m|}\bigl(g_m-\kappa_m+1\bigr)
-\frac{1}{|X_m|}\sum_{w\in X_m}\log d^{\mathrm{bin}}_m(w).
\]
则 \(\Delta_m=\log C_m\)。因而对任意固定整数 \(R\ge 1\)，当 \(m\to\infty\) 时有
\[
\Delta_m
=
\log(2\pi)+
\sum_{r=1}^{R}
\frac{B_{2r}}{r(2r-1)}
\Bigl(\varphi^{-1}+\varphi^{2r-3}\Bigr)
\Bigl(\frac{\varphi}{2}\Bigr)^{(2r-1)m}
+O\!\Bigl(\Bigl(\frac{\varphi}{2}\Bigr)^{(2R+1)m}\Bigr).
\]
\end{proposition}
\begin{proof}
由推论 \ref{cor:fold-bin-recover-2pi} 中 \(C_m\) 的定义，立即有
\[
\Delta_m=\log C_m.
\]
再将定理 \ref{thm:fold-bin-gauge-constant-stirling-bernoulli-hierarchy} 对 \(\log C_m\) 的完整渐近展开代入，即得所示公式。
\end{proof}

\begin{corollary}[归一化规范亏损尾部的偶 \texorpdfstring{\(\zeta\)}{zeta} 刚性]\label{cor:killo-fold-bin-normalized-gauge-deficiency-tail-rigidity}
\leanverified{paper\_killo\_fold\_bin\_normalized\_gauge\_deficiency\_tail\_rigidity}
设 \((\Delta_m^{(1)})_{m\ge 1}\) 与 \((\Delta_m^{(2)})_{m\ge 1}\) 为两列满足命题 \ref{prop:killo-fold-bin-normalized-gauge-deficiency} 渐近律的归一化规范亏损序列。
若存在 \(m_0\) 使得
\[
\Delta_m^{(1)}=\Delta_m^{(2)}
\qquad (m\ge m_0),
\]
则由这两列序列恢复出的 Bernoulli 系数与偶 \(\zeta\)-谱逐项一致。
更确切地，若记对应恢复值为 \(B_{2r}^{(j)}\) 与 \(\zeta^{(j)}(2r)\)（\(j\in\{1,2\}\)），则对一切 \(r\ge 1\) 有
\[
B_{2r}^{(1)}=B_{2r}^{(2)},
\qquad
\zeta^{(1)}(2r)=\zeta^{(2)}(2r).
\]
\end{corollary}
\begin{proof}
由命题 \ref{prop:killo-fold-bin-normalized-gauge-deficiency}，两列序列分别满足与 \(\log C_m\) 相同的指数层级展开。
而命题 \ref{prop:fold-bin-gauge-bernoulli-extraction-operator} 给出的抽取算子只依赖于充分大的尾部值；因此一旦 \(\Delta_m^{(1)}\) 与 \(\Delta_m^{(2)}\) 在所有充分大 \(m\) 上相同，逐阶抽取得到的全部 \(B_{2r}\) 必然相同。
最后由定理 \ref{thm:fold-bin-gauge-constant-zeta-even-recovery-limit} 中 Bernoulli 数与偶 \(\zeta\)-值的显式对应，得到所述偶谱一致性。
\end{proof}

\endinput


\begin{remark}[两态压缩作为解析基准解]\label{rem:fold-bin-solvable-model}
在 bin-fold 情形下，命题 \ref{prop:fold-bin-power-sum-asymptotic}--命题 \ref{prop:fold-bin-gauge-density-two-scale} 表明：全阶碰撞矩、全温度 escort 族、Rényi/\(f\)-散度族以及规范体积密度的首个可见修正项，均可在同一对两态常数与一致误差阶下闭式化。
该结构为正文更一般折叠机制提供一个可对照的解析基准。
\end{remark}

\begin{theorem}[bin-fold 的尾谱极限诱导的线性预言机阈值与指数律]\label{thm:fold-bin-oracle-linear-threshold-exponent}
\leanverified{paper\_fold\_bin\_oracle\_linear\_threshold\_exponent}
在 bin-fold 口径下记 \(d^{\mathrm{bin}}_m(w):=\bigl|(\mathrm{Fold}^{\mathrm{bin}}_m)^{-1}(w)\bigr|\)（\(w\in X_m\)）。
对整数 \(B\ge 0\) 定义预言机容量与最优成功率
\[
\mathcal C^{\mathrm{bin}}_m(B):=\sum_{w\in X_m}\min\bigl(d^{\mathrm{bin}}_m(w),2^B\bigr),
\qquad
\mathrm{Succ}^{\mathrm{bin}}_m(B):=\frac{\mathcal C^{\mathrm{bin}}_m(B)}{2^m}.
\]
令
\[
\alpha_c:=\log_2\!\Bigl(\frac{2}{\varphi}\Bigr)=1-\log_2\varphi.
\]
取 \(B_m=\lfloor \alpha m\rfloor\)。
则：
\begin{enumerate}
  \item 若 \(\alpha>\alpha_c\)，则存在 \(m_0\) 使对所有 \(m\ge m_0\) 有 \(\mathrm{Succ}^{\mathrm{bin}}_m(B_m)=1\)。
  \item 若 \(\alpha<\alpha_c\)，则极限存在且满足
  \[
  \lim_{m\to\infty}\frac{1}{m}\log_2 \mathrm{Succ}^{\mathrm{bin}}_m(B_m)=\alpha-\alpha_c.
  \]
\end{enumerate}
\end{theorem}
\begin{proof}
由定义 \(\mathcal C_m^{\mathrm{bin}}(B)=\mathcal C_m^{\mathrm{bin,cont}}(2^B)\)，其中 \(\mathcal C_m^{\mathrm{bin,cont}}\) 为命题 \ref{prop:fold-bin-degeneracy-tail-capacity-kinks} 的连续化容量曲线。
并且 \(\mathrm{Succ}^{\mathrm{bin}}_m(B)=2^{-m}\mathcal C_m^{\mathrm{bin,cont}}(2^B)\)。
\par
令 \(D^{\mathrm{bin}}_m:=\max_{w\in X_m}d_m^{\mathrm{bin}}(w)\)。
由定理 \ref{thm:fold-bin-two-state-asymptotic} 的一致展开可知
\(
D^{\mathrm{bin}}_m=(2/\varphi)^m+O(1)
\)，从而 \(D^{\mathrm{bin}}_m\le 2^{\alpha m}\) 对充分大 \(m\) 成立（当 \(\alpha>\alpha_c\)）。
于是对所有 \(w\) 有 \(\min(d_m^{\mathrm{bin}}(w),2^{B_m})=d_m^{\mathrm{bin}}(w)\)，从而
\(
\mathcal C_m^{\mathrm{bin}}(B_m)=\sum_w d_m^{\mathrm{bin}}(w)=2^m
\)，即 \(\mathrm{Succ}^{\mathrm{bin}}_m(B_m)=1\)。
\par
当 \(\alpha<\alpha_c\) 时，对充分大 \(m\) 有 \(2^{B_m}\le \varphi^{-1}(2/\varphi)^m+O(1)\)，从而 \(2^{B_m}\le \min_w d_m^{\mathrm{bin}}(w)\) 成立。
于是逐点有 \(\min(d_m^{\mathrm{bin}}(w),2^{B_m})=2^{B_m}\)，得到
\[
\mathcal C_m^{\mathrm{bin}}(B_m)=|X_m|\,2^{B_m},
\qquad
\mathrm{Succ}^{\mathrm{bin}}_m(B_m)=\frac{|X_m|}{2^m}\,2^{B_m}.
\]
由 \(|X_m|=F_{m+2}=\varphi^m\cdot \Theta(1)\) 得
\(
\frac{1}{m}\log_2(|X_m|/2^m)\to \log_2(\varphi/2)=-\alpha_c
\)，
又 \(B_m/m\to \alpha\)，合并即得所示指数律。证毕。
\end{proof}

\begin{theorem}[临界舍入的 log-periodic 非收敛与全子列极限区间]\label{thm:fold-bin-critical-rounding-logperiodic}
\leanverified{paper\_fold\_bin\_critical\_rounding\_logperiodic}
令 \(\alpha_c=\log_2(2/\varphi)\) 且 \(B_m=\lfloor \alpha_c m\rfloor\)。
定义
\[
r_m:=2^{B_m-\alpha_c m}\in(1/2,1].
\]
则 \(\mathrm{Succ}^{\mathrm{bin}}_m(B_m)\) 不收敛；其全部子列极限集合恰为闭区间
\[
\biggl[\frac{\varphi^2}{2\sqrt5},\,1\biggr].
\]
更精确地，若沿子列满足 \(r_m\to r\in[1/2,1]\)，则
\[
\lim \mathrm{Succ}^{\mathrm{bin}}_m(B_m)
=
\begin{cases}
\dfrac{\varphi^2}{\sqrt5}\,r, & 1/2\le r\le\varphi^{-1},\\[10pt]
\dfrac{r+\varphi^{-2}}{1+\varphi^{-2}}, & \varphi^{-1}\le r\le 1.
\end{cases}
\]
\end{theorem}
\begin{proof}
先证 \(\alpha_c\notin\QQ\)。
若 \(\alpha_c=p/q\in\QQ\)（互素 \(p,q\)），则 \(2^{\alpha_c}=2^{p/q}\) 为代数数且满足
\(
2^{\alpha_c}=2/\varphi
\)，从而 \(\varphi^q=2^{q-p}\in\QQ\)。
但 \(\varphi\) 为二次无理数，且对任意 \(q\ge 1\) 都有 \(\varphi^q\notin\QQ\)，矛盾。
\par
由 \(\alpha_c\notin\QQ\) 与 Kronecker 定理，\(\{m\alpha_c\}\) 在 \([0,1]\) 稠密，从而
\(
r_m=2^{-\{m\alpha_c\}}
\)
在 \((1/2,1]\) 稠密。
\par
另一方面，由命题 \ref{prop:fold-bin-degeneracy-tail-capacity-kinks} 的尺度化容量过程 \(\widehat C_m^{\mathrm{bin}}(s)\)（其中
\(\widehat C_m^{\mathrm{bin}}(s)=2^{-m}\mathcal C_m^{\mathrm{bin,cont}}(s(2/\varphi)^m)\)），
并注意 \(2^{B_m}=r_m(2/\varphi)^m\)，得
\[
\mathrm{Succ}^{\mathrm{bin}}_m(B_m)
=2^{-m}\mathcal C_m^{\mathrm{bin,cont}}(2^{B_m})
=\widehat C_m^{\mathrm{bin}}(r_m).
\]
由该命题，\(\widehat C_m^{\mathrm{bin}}\) 在 \([1/2,1]\) 上一致收敛到连续极限函数
\[
\widehat C_\infty^{\mathrm{bin}}(s)
=\frac{\varphi}{\sqrt5}\min\{1,s\}+\frac{1}{\sqrt5}\min\{\varphi^{-1},s\}.
\]
因此若 \(r_m\to r\)，则 \(\mathrm{Succ}^{\mathrm{bin}}_m(B_m)\to \widehat C_\infty^{\mathrm{bin}}(r)\)。
将 \(\widehat C_\infty^{\mathrm{bin}}\) 在区间 \([1/2,1]\) 上逐段展开即得所示分段闭式，
并且 \(\widehat C_\infty^{\mathrm{bin}}([1/2,1])=[\widehat C_\infty^{\mathrm{bin}}(1/2),\widehat C_\infty^{\mathrm{bin}}(1)]=[\varphi^2/(2\sqrt5),1]\)。
由于 \(r_m\) 稠密，子列极限集合即为该闭区间，且整体序列不收敛。证毕。
\end{proof}

\begin{theorem}[bin-fold 输出分布的 Rényi 质量谱仿射律]\label{thm:fold-bin-renyi-mass-spectrum-affine}
\leanverified{paper\_fold\_bin\_renyi\_mass\_spectrum\_affine}
令 \(\pi_m(w):=d_m^{\mathrm{bin}}(w)/2^m\) 为 bin-fold 在均匀微态输入下诱导的输出分布。
则对任意 \(q>0\)，存在常数 \(K_q>0\) 使
\[
\sum_{w\in X_m}\pi_m(w)^q
=
K_q\cdot \varphi^{(1-q)m}\cdot\bigl(1+o(1)\bigr),
\qquad
K_q=\frac{\varphi+\varphi^{-q}}{\sqrt5}.
\]
从而 Rényi 质量谱（自然对数）满足
\[
h_q:=-\lim_{m\to\infty}\frac{1}{m}\log\sum_{w\in X_m}\pi_m(w)^q
=(q-1)\log\varphi.
\]
\end{theorem}
\begin{proof}
由 \(\pi_m(w)^q=2^{-qm}\bigl(d_m^{\mathrm{bin}}(w)\bigr)^q\)，有
\[
\sum_{w\in X_m}\pi_m(w)^q
=2^{-qm}S_q^{\mathrm{bin}}(m).
\]
由命题 \ref{prop:fold-bin-power-sum-asymptotic}，
\[
S_q^{\mathrm{bin}}(m)
=\Bigl(\frac{2}{\varphi}\Bigr)^{qm}\bigl(F_{m+1}+\varphi^{-q}F_m\bigr)\bigl(1+o(1)\bigr),
\]
从而
\[
\sum_{w\in X_m}\pi_m(w)^q
=\varphi^{-qm}\bigl(F_{m+1}+\varphi^{-q}F_m\bigr)\bigl(1+o(1)\bigr).
\]
由 Binet 渐近 \(F_{m+1}=\varphi^{m+1}/\sqrt5+o(\varphi^m)\)、\(F_m=\varphi^m/\sqrt5+o(\varphi^m)\)，得
\[
\sum_{w\in X_m}\pi_m(w)^q
=\frac{\varphi^{m+1-qm}+\varphi^{m-qm-q}}{\sqrt5}\bigl(1+o(1)\bigr)
=\frac{\varphi+\varphi^{-q}}{\sqrt5}\,\varphi^{(1-q)m}\bigl(1+o(1)\bigr),
\]
即所示常数 \(K_q\) 与主阶。
取 \(-\frac{1}{m}\log\) 并令 \(m\to\infty\) 得 \(h_q=(q-1)\log\varphi\)。证毕。
\end{proof}

\endinput



\begin{corollary}[重建预算的相变：阈值 \texorpdfstring{$\varepsilon_c=\varphi/\sqrt5$}{epsc=phi/sqrt5} 处出现 \texorpdfstring{$\log\varphi$}{log phi} 级节省]\label{cor:fold-bin-budget-phase-transition}
\leanverified{paper\_fold\_bin\_budget\_phase\_transition}
沿用推论 \ref{cor:fold-bin-quantile-threshold-constant} 中的尾部分位阈值 \(B^{\mathrm{bin}}_m(\varepsilon)\)。
若某类重建/解码协议的失败概率可由纤维阈值控制为
\[
\mathbb{P}\Bigl(d^{\mathrm{bin}}_m(W_m)>\frac{P_m}{2}\Bigr)\ \le\ \varepsilon,
\]
则充分条件是选择模积预算 \(P_m>2B^{\mathrm{bin}}_m(\varepsilon)\)。
在该口径下，令 \(\varepsilon_c:=\varphi/\sqrt5\)。
则最小对数预算呈现分段相变：
\[
\log P_m\ \ge\ m\log\!\Bigl(\frac{2}{\varphi}\Bigr)+\log 2+o(1)\qquad(0<\varepsilon<\varepsilon_c),
\]
而当 \(\varepsilon>\varepsilon_c\) 时可降低为
\[
\log P_m\ \ge\ m\log\!\Bigl(\frac{2}{\varphi}\Bigr)+\log 2-\log\varphi+o(1).
\]
\end{corollary}

\begin{proof}
由 \(B^{\mathrm{bin}}_m(\varepsilon)\) 的定义，恒有
\(
\mathbb{P}\bigl(d^{\mathrm{bin}}_m(W_m)>B^{\mathrm{bin}}_m(\varepsilon)\bigr)\le \varepsilon
\)；
因此若 \(P_m>2B^{\mathrm{bin}}_m(\varepsilon)\) 则
\(
\mathbb{P}\bigl(d^{\mathrm{bin}}_m(W_m)>P_m/2\bigr)\le\varepsilon
\)，满足所需的失败概率控制。
\par
另一方面，由推论 \ref{cor:fold-bin-quantile-threshold-constant}，
\[
\Bigl(\frac{\varphi}{2}\Bigr)^m B^{\mathrm{bin}}_m(\varepsilon)\to
\begin{cases}
1,& 0<\varepsilon<\varepsilon_c,\\
\varphi^{-1},& \varepsilon_c<\varepsilon<1,
\end{cases}
\]
从而
\[
\log B^{\mathrm{bin}}_m(\varepsilon)
=m\log\!\Bigl(\frac{2}{\varphi}\Bigr)+
\begin{cases}
o(1),& 0<\varepsilon<\varepsilon_c,\\
-\log\varphi+o(1),& \varepsilon>\varepsilon_c.
\end{cases}
\]
取 \(P_m=2B^{\mathrm{bin}}_m(\varepsilon)(1+o(1))\) 即得两条下界。证毕。
\end{proof}

\begin{corollary}[最大纤维、组合上界与饱和判别]\label{cor:fold-bin-saturation}
\leanverified{paper\_fold\_bin\_saturation}
对任意 \(m\ge 1\)，最大纤维由 \(w=0^m\) 取得，并满足组合上界
\begin{equation}\label{eq:fold-bin-upper}
d^{\mathrm{bin}}_{\max}(m):=d^{\mathrm{bin}}_m(0^m)\le |X_{K(m)-m}|=F_{K(m)-m+2}.
\end{equation}
记 \(K=K(m)\)、\(L=K-m\)，并定义在相邻约束下的\emph{最大尾和}
\[
S_{\max}(m):=\max\Bigl\{\sum_{k=m+1}^{K} c_k F_{k+1}:\ (c_{m+1},\dots,c_K)\in\{0,1\}^{L},\ c_k c_{k+1}=0\Bigr\}.
\]
则对 \(w=0^m\)，上界 \eqref{eq:fold-bin-upper} 饱和当且仅当阈值约束对所有合法尾部都是冗余的，即
\begin{equation}\label{eq:fold-bin-sat-iff}
d^{\mathrm{bin}}_{\max}(m)=F_{K(m)-m+2}\quad\Longleftrightarrow\quad 2^m-1\ge S_{\max}(m).
\end{equation}
并且 \(S_{\max}(m)\) 具有完全显式的闭式（由“1010\(\cdots\)”极大尾部取得）：
\begin{equation}\label{eq:fold-bin-Smax-closed}
S_{\max}(m)=
\begin{cases}
F_{K(m)+2}-F_{m+1}, & L=K(m)-m\ \text{为奇数},\\
F_{K(m)+2}-F_{m+2}, & L=K(m)-m\ \text{为偶数}.
\end{cases}
\end{equation}
等价地，饱和判别可改写为
\[
F_{K(m)+2}-2^m\le
\begin{cases}
F_{m+1}-1,& L\ \text{奇},\\
F_{m+2}-1,& L\ \text{偶}.
\end{cases}
\]
\end{corollary}

\begin{corollary}[饱和判别的余数窗口形式：顶端区间条件]\label{cor:fold-bin-saturation-residue-window}
沿用推论 \ref{cor:fold-bin-saturation} 的记号，令 \(K=K(m)\) 并置
\[
r_m:=(2^m-1)-F_{K+1}\ \in\ \{0,1,\dots,F_K-1\}.
\]
则推论 \ref{cor:fold-bin-saturation} 的饱和条件
\(
2^m-1\ge S_{\max}(m)
\)
等价于如下“顶端窗口”条件：
\[
r_m\in\bigl[F_K-F_{m+1},\,F_K-1\bigr]\quad (L\ \text{奇}),
\qquad
r_m\in\bigl[F_K-F_{m+2},\,F_K-1\bigr]\quad (L\ \text{偶}).
\]
因此饱和并非由 \(2^m\) 的规模决定，而由其相对 Fibonacci 阶梯 \(F_{K+1}\) 的余数落点所决定。
\end{corollary}
\leanpartial{paper\_fold\_bin\_saturation\_residue\_window}{the current repository seeds define `foldBinDigitK m = 2 * m`, which makes the stated residue window false already at `m = 3`}
\begin{proof}
由 \eqref{eq:fold-bin-Smax-closed} 与 \(F_{K+2}=F_{K+1}+F_K\)，
不等式 \(2^m-1\ge S_{\max}(m)\) 等价于
\[
(2^m-1)-F_{K+1}\ \ge\
\begin{cases}
F_K-F_{m+1},& L\ \text{奇},\\
F_K-F_{m+2},& L\ \text{偶},
\end{cases}
\]
即所示余数窗口条件。证毕。
\end{proof}

\begin{corollary}[饱和窗口的相对尺度与指数衰减]\label{cor:fold-bin-saturation-window-density}
在推论 \ref{cor:fold-bin-saturation-residue-window} 的余数表示下，饱和条件仅要求 \(r_m\) 落入长度
\[
\ell_m=
\begin{cases}
F_{m+1},& L\ \text{奇},\\
F_{m+2},& L\ \text{偶}
\end{cases}
\]
的顶端窗口中；相对于全余数区间长度 \(F_K\)，其占比为 \(\ell_m/F_K\)。
由于 \(K(m)=m\log_{\varphi}2+O(1)\)，该占比满足
\[
\frac{\ell_m}{F_{K(m)}}=\Theta\!\bigl(\varphi^{m-K(m)}\bigr)=\Theta\!\bigl(\varphi^{-(K(m)-m)}\bigr),
\]
从而随 \(m\) 指数级衰减；该“顶端窗口”图景与注 \ref{rem:fold-bin-finite} 的对数线性形式刚性在机制层面一致。
\end{corollary}

\begin{theorem}[最大纤维指数与 \texorpdfstring{$\log(2/\varphi)$}{log(2/phi)}（bin-fold）]\label{thm:fold-bin-max-fiber-exponent}
令 \(d^{\mathrm{bin}}_{\max}(m):=\max_{w\in X_m} d^{\mathrm{bin}}_m(w)\)。
则极限存在并满足
\[
\lim_{m\to\infty}\frac{1}{m}\log d^{\mathrm{bin}}_{\max}(m)=\log\frac{2}{\varphi}.
\]
特别地，该指数常数与正文命题 \ref{prop:beta-quantization-step} 中的不可约步长 \(\Delta_{\min}=\log(2/\varphi)\) 同一。
\leanverified{paper\_fold\_bin\_max\_fiber\_exponent}
\end{theorem}
\begin{proof}
由满射 \(\mathrm{Fold}^{\mathrm{bin}}_m:\{0,\dots,2^m-1\}\twoheadrightarrow X_m\) 的计数恒等式
\(
\sum_{w\in X_m} d^{\mathrm{bin}}_m(w)=2^m
\)
得
\(
d^{\mathrm{bin}}_{\max}(m)\ge 2^m/|X_m|
=2^m/F_{m+2}
\)。
由 Binet 渐近 \(F_{m+2}=\Theta(\varphi^{m})\) 得
\[
\liminf_{m\to\infty}\frac{1}{m}\log d^{\mathrm{bin}}_{\max}(m)\ \ge\ \log\frac{2}{\varphi}.
\]
\par
另一方面，由推论 \ref{cor:fold-bin-saturation} 的上界
\(
d^{\mathrm{bin}}_{\max}(m)\le F_{K(m)-m+2}
\)。
由定义 \eqref{eq:K-of-m} 有
\(
F_{K(m)+1}\le 2^m-1 < F_{K(m)+2}
\)，
结合 \(\log F_n=n\log\varphi+O(1)\) 得
\[
\bigl(K(m)+1\bigr)\log\varphi+O(1)\ \le\ m\log 2\ \le\ \bigl(K(m)+2\bigr)\log\varphi+O(1),
\]
从而 \(K(m)=m\log_{\varphi}2+O(1)\)，故
\[
\log F_{K(m)-m+2}=\bigl(K(m)-m\bigr)\log\varphi+O(1)
=m(\log 2-\log\varphi)+O(1).
\]
两边同除以 \(m\) 并取上极限，得
\(
\limsup_{m\to\infty}\frac{1}{m}\log d^{\mathrm{bin}}_{\max}(m)\le \log(2/\varphi)
\)。
合并上下界即得极限同一性。证毕。
\end{proof}

\begin{remark}[不可约步长常数的硬约束]\label{rem:fold-bin-delta-min-hard-constraint}
定理 \ref{thm:fold-bin-max-fiber-exponent} 给出：在 bin-fold 的纤维结构不变的前提下，最大可见冗余增长率的指数常数被强制为 \(\log(2/\varphi)\)。
因此，若某一方案在同一折叠纪律下声称实现严格小于 \(\log(2/\varphi)\) 的不可约步长（正文命题 \ref{prop:beta-quantization-step} 的意义），则与最大纤维指数定理所提供的增长率证书发生结构性冲突。
\end{remark}

\begin{remark}[指数级 Diophantine 近似与“仅有限次饱和”]\label{rem:fold-bin-finite}
由 \eqref{eq:fold-bin-Smax-closed} 可见，饱和条件强迫 \(2^m\) 与某个 Fibonacci 数 \(F_{K(m)+2}\) \emph{异常接近}。
用 Binet 公式 \(F_n=\varphi^n/\sqrt5+O(\varphi^{-n})\) 将其改写，可导出一个指数级小的线性对数形式约束：
\[
\Bigl|(K(m)+2)\log\varphi - m\log 2 - \log\sqrt5\Bigr|
\ \lesssim\ \Bigl(\frac{\varphi}{2}\Bigr)^m.
\]
由于 \(\varphi\) 为代数数而 \(2\) 为有理数，上式左侧是非零的对数线性形式；而 Baker 型下界（例如 Matveev 的显式形式）给出其\emph{至多多项式级}的小下界（见 \cite{Matveev2000LinearFormsLogs,BakerWustholz2007LogarithmicForms}），从而与右侧的指数级小量矛盾，推出：
\[
\text{上界饱和只能发生有限次。}
\]
本结论在本文的可复现管线中可直接数值审计（见表 \ref{tab:fold_binary_saturation_points}）。
\end{remark}

\leanverified{paper\_fold\_bin\_digit\_dp}
\begin{proposition}[计算 \(\boldsymbol{d^{\mathrm{bin}}_m(w)}\) 的 \(O(m)\) Fibonacci-digit DP]\label{prop:fold-bin-digit-dp}
给定 \(m\ge 1\) 与 \(w\in X_m\)。令 \(K=K(m)\) 并记 \(V=V_m(w)\)，再令 \(\mathrm{slack}:=2^m-1-V\)。
将 \(\mathrm{slack}\) 写为 Zeckendorf 位串 \(Z(\mathrm{slack})=(b_1,\dots,b_K)\)（合法，无相邻 \(11\)）。
则 \(d^{\mathrm{bin}}_m(w)\) 可由如下 digit DP 在 \(O(K-m)=O(m)\) 时间内精确计算：
\[
\mathrm{DP}(k,\mathrm{prev},\mathrm{tight})
\quad(k=K,K-1,\dots,m+1),
\]
其中 \(\mathrm{prev}\in\{0,1\}\) 记录上一个（更高位）是否取 \(1\)，\(\mathrm{tight}\in\{0,1\}\) 表示当前前缀是否仍与 \((b_k,\dots,b_K)\) 相同；
转移时允许选择 \(c_k\in\{0,1\}\)，满足
\[
c_k c_{k+1}=0,\qquad
\mathrm{tight}=1\Rightarrow c_k\le b_k,\qquad
w_m=1\Rightarrow c_{m+1}=0,
\]
并按是否保持等号更新 \(\mathrm{tight}\)。该 DP 的输出即为满足 \(\sum_{k=m+1}^K c_k F_{k+1}\le \mathrm{slack}\) 的尾部数目，从而等于 \(d^{\mathrm{bin}}_m(w)\)。
\end{proposition}

\begin{table}[H]
\centering
\scriptsize
\setlength{\tabcolsep}{6pt}
\caption{二进制区间折叠 $\mathrm{Fold}^{\mathrm{bin}}_m$ 的“上界饱和点”（扫描 $m\le 400$）。表中给出 $K(m)$（满足 $F_{K(m)+1}\le 2^m-1 < F_{K(m)+2}$）、尾长 $L=K-m$、组合上界 $F_{L+2}$、精确最大纤维 $d_{\max}(m)=|\mathrm{Fold}^{\mathrm{bin}}_m{}^{-1}(0^m)|$（digit DP）、以及最大尾和 $S_{\max}(m)$ 与阈值 $2^m-1$ 的裕量。}
\label{tab:fold_binary_saturation_points}
\begin{tabular}{rrrrrrrr}
\toprule
$m$ & $K(m)$ & $L$ & $F_{L+2}$ & $d_{\max}(m)$ & $S_{\max}(m)$ & $2^m-1$ & $(2^m-1)-S_{\max}$\\
\midrule
1 & 1 & 0 & 1 & 1 & 0 & 1 & 1\\
2 & 3 & 1 & 2 & 2 & 3 & 3 & 0\\
3 & 4 & 1 & 2 & 2 & 5 & 7 & 2\\
4 & 6 & 2 & 3 & 3 & 13 & 15 & 2\\
5 & 7 & 2 & 3 & 3 & 21 & 31 & 10\\
7 & 10 & 3 & 5 & 5 & 123 & 127 & 4\\
12 & 17 & 5 & 13 & 13 & 3948 & 4095 & 147\\
\bottomrule
\end{tabular}
\vspace{4pt}
\\饱和点集合（$m\le 400$）：$\{ 1,\,2,\,3,\,4,\,5,\,7,\,12 \}$.
\end{table}

